%--------------------------------------------------------------------%
%
% Berkas utama templat Skripsi Universitas Negeri Malang.
%
% Diadaptasi dari templat Tesis IF-ITB (Petra Barus, Peb Ruswono Aryan)
% mengikuti Pedoman Penulisan Karya Ilmiah UM Edisi 2017.
%
%--------------------------------------------------------------------%

\documentclass[12pt, a4paper, onecolumn, oneside, final]{report}

\input{config/if-itb-thesis.sty}
%-------------------------------------------------------------------%
%
% Informasi umum skripsi — SUNTING DI SINI SAJA.
% Semua halaman front matter mengambil nilai dari makro di bawah.
%
%-------------------------------------------------------------------%

\newcommand{\judul}{Weight-Only Post-Training Quantization untuk Extractive Question Answering Bahasa Indonesia}
\newcommand{\penulis}{Kardynan Parulian}
\newcommand{\nim}{230535607270}

\newcommand{\departemen}{Teknik Elektro dan Informatika}
\newcommand{\fakultas}{Fakultas Teknik}
\newcommand{\programstudi}{Program Studi Teknik Informatika}
% Diisi pada halaman judul: "...menyelesaikan program <gelarprogram>"
\newcommand{\gelarprogram}{Sarjana Teknik Informatika}

\newcommand{\bulantahun}{Oktober 2026}
\newcommand{\kotatanggal}{Malang, 23 Oktober 2026}

\newcommand{\pembimbingSatu}{Muhammad Jauharul Fuady, S.T., M.T.}
\newcommand{\nipSatu}{NIP 198305072009121002}
% Pembimbing II: kosongkan (isi kurung kurawal) bila hanya ada satu pembimbing
\newcommand{\pembimbingDua}{}
\newcommand{\nipDua}{}

%-------------------------------------------------------------------%
% Lembar Persetujuan dan Pengesahan (Lampiran 5a)
% Periksa kembali nama pejabat saat berkas dicetak.
%-------------------------------------------------------------------%

% Tanggal sidang; kosongkan isinya bila belum ditetapkan (tercetak garis kosong)
\newcommand{\tanggalSidang}{\underline{\hspace{3cm}}}

% Dewan penguji: kosongkan bila belum ditetapkan. Baris nama dan NIP
% akan tercetak sebagai garis kosong untuk diisi tangan.
\newcommand{\pengujiSatu}{}
\newcommand{\nipPengujiSatu}{}
\newcommand{\pengujiDua}{}
\newcommand{\nipPengujiDua}{}
\newcommand{\pengujiTiga}{}
\newcommand{\nipPengujiTiga}{}

% Pejabat penanda tangan
\newcommand{\dekan}{Prof. Dr. Andoko, S.T., M.T.}
\newcommand{\nipDekan}{NIP 196508121991031005}
\newcommand{\ketuaDepartemen}{Dr. Ir. Triyanna Widiyaningtyas, M.T.}
\newcommand{\nipKetuaDepartemen}{NIP 197412152008122002}

% Kata kunci ringkasan (3--5 kata), dipisah koma
\newcommand{\katakunci}{kuantisasi, weight-only, granularitas kuantisasi, extractive question answering, IndoBERT}
\newcommand{\katakunciEn}{post-training quantization, weight-only, quantization granularity, extractive question answering, IndoBERT}


\bibliography{references}

\begin{document}

    \title{\judul}
    \author{\penulis}
    \date{}

    %----------------------------------------------------------------%
    % Bagian awal — nomor halaman romawi, tengah bawah.
    %----------------------------------------------------------------%
    \pagenumbering{roman}
    \setcounter{page}{1}

    % Halaman sampul, lambang, dan halaman judul tidak bernomor dan tidak
    % dihitung, sehingga Lembar Persetujuan menjadi halaman i.
    %--------------------------------------------------------------------%
%
% Halaman Sampul Skripsi (Pedoman UM Lampiran 1a)
%
% Nilai diambil dari config/informations.tex.
%
%--------------------------------------------------------------------%

\clearpage
\thispagestyle{empty}

\begin{center}
    \setstretch{1.5}

    \bfseries \MakeUppercase{\judul}
    \vfill

    SKRIPSI
    \vfill

    OLEH \\
    \MakeUppercase{\penulis} \\
    NIM \nim%
    \vfill

    % Lambang UM (diameter 3 cm)
    \includegraphics[width=3cm]{resources/logo-um.png}
    \vfill

    \MakeUppercase{Universitas Negeri Malang} \\
    \MakeUppercase{\fakultas} \\
    \MakeUppercase{\programstudi} \\
    \MakeUppercase{\bulantahun}

\end{center}

\clearpage

    %--------------------------------------------------------------------%
%
% Lembar Logo (Pedoman UM Lampiran 2) — hanya lambang UM diameter 8 cm.
%
%--------------------------------------------------------------------%

\clearpage
\thispagestyle{empty}

\begin{center}
    \vspace*{\fill}
    % Lambang UM (diameter 8 cm)
    \includegraphics[width=8cm]{resources/logo-um.png}
    \vspace*{\fill}
\end{center}

\clearpage

    %--------------------------------------------------------------------%
%
% Halaman Judul lembar kedua (Pedoman UM Lampiran 3a)
%
%--------------------------------------------------------------------%

\clearpage
\thispagestyle{empty}

\begin{center}
    \setstretch{1.5}

    \bfseries \MakeUppercase{\judul}
    \vfill

    \normalfont%
    SKRIPSI \\[0.3\baselineskip]
    Diajukan kepada \\
    Universitas Negeri Malang \\
    untuk memenuhi salah satu persyaratan \\
    dalam menyelesaikan program \gelarprogram%
    \vfill

    OLEH \\
    \MakeUppercase{\penulis} \\
    NIM \nim%
    \vfill

    % Lambang UM (diameter 3 cm)
    \includegraphics[width=3cm]{resources/logo-um.png}
    \vfill

    \bfseries
    \MakeUppercase{Universitas Negeri Malang} \\
    \MakeUppercase{\fakultas} \\
    \MakeUppercase{\programstudi} \\
    \MakeUppercase{\bulantahun}

\end{center}

\clearpage


    \setcounter{page}{1}
    \pagestyle{plain}

    %--------------------------------------------------------------------%
%
% Lembar Persetujuan Pembimbing (Pedoman UM Lampiran 4a)
%
% Halaman ini adalah halaman pertama bagian awal yang bernomor (i).
%
%--------------------------------------------------------------------%

\clearpage

\begin{center}
    \bfseries LEMBAR PERSETUJUAN PEMBIMBING SKRIPSI
\end{center}

\phantomsection
\addcontentsline{toc}{chapter}{LEMBAR PERSETUJUAN PEMBIMBING SKRIPSI}

\vspace{\baselineskip}

\noindent Skripsi oleh \penulis\ dengan judul ``\judul'' telah diperiksa dan disetujui untuk dilaksanakan sidang skripsi.

\vspace{2\baselineskip}

% Satu pembimbing: satu blok tanda tangan rata kiri.
% Dua pembimbing: dua blok berdampingan (Pembimbing I dan II).
\ifx\pembimbingDua\empty
\noindent
\begin{tabular}{@{}l@{}}
\kotatanggal \\
Dosen Pembimbing \\[2cm]
\pembimbingSatu \\
\nipSatu
\end{tabular}
\else
\noindent
\begin{tabular}{@{}p{0.45\textwidth}p{0.45\textwidth}@{}}
\kotatanggal & \kotatanggal \\
Pembimbing I & Pembimbing II \\[2cm]
\pembimbingSatu & \pembimbingDua \\
\nipSatu & \nipDua
\end{tabular}
\fi

\clearpage
      % Lembar Persetujuan Pembimbing
    %--------------------------------------------------------------------%
%
% Lembar Persetujuan dan Pengesahan (Pedoman UM Lampiran 5a)
%
% Nama dewan penguji dan pejabat penanda tangan diambil dari
% config/informations.tex.
%
%--------------------------------------------------------------------%

\clearpage

\begin{center}
    \bfseries LEMBAR PERSETUJUAN DAN PENGESAHAN SKRIPSI
\end{center}

\phantomsection
\addcontentsline{toc}{chapter}{LEMBAR PERSETUJUAN DAN PENGESAHAN SKRIPSI}

\vspace{\baselineskip}

\noindent Skripsi oleh \MakeUppercase{\penulis} ini telah dipertahankan di depan dewan penguji pada tanggal \tanggalSidang.

\vspace{\baselineskip}

\noindent Dewan Penguji

\vspace{\baselineskip}

% Satu baris penguji: nama (atau garis kosong) lalu NIP (atau garis kosong).
\newcommand*{\pengujiBaris}[2]{%
  \ifx#1\empty\underline{\strut\hspace{6.5cm}}\else#1\fi\\[0.15cm]
  \ifx#2\empty NIP~\underline{\strut\hspace{4.5cm}}\else#2\fi}

\noindent
\pengujiBaris{\pengujiSatu}{\nipPengujiSatu}

\vspace{0.8cm}
\noindent
\pengujiBaris{\pengujiDua}{\nipPengujiDua}

\vspace{0.8cm}
\noindent
\pengujiBaris{\pengujiTiga}{\nipPengujiTiga}

\vspace{2\baselineskip}

\noindent
\begin{tabular}{@{}p{0.48\textwidth}p{0.48\textwidth}@{}}
Mengesahkan, & Mengetahui, \\
Dekan \fakultas, & Ketua Departemen \departemen, \\[2cm]
\dekan & \ketuaDepartemen \\
\nipDekan & \nipKetuaDepartemen
\end{tabular}

\clearpage
   % Lembar Persetujuan dan Pengesahan
    %--------------------------------------------------------------------%
%
% Pernyataan Keaslian Tulisan (Pedoman UM Lampiran 6)
%
%--------------------------------------------------------------------%

\clearpage

\begin{center}
    \bfseries PERNYATAAN KEASLIAN TULISAN
\end{center}

\phantomsection
\addcontentsline{toc}{chapter}{PERNYATAAN KEASLIAN TULISAN}

\vspace{\baselineskip}
\noindent Saya yang bertanda tangan di bawah ini:

\vspace{0.5\baselineskip}
\noindent
\begin{tabular}{@{}p{5cm}@{: }>{\raggedright\arraybackslash}p{8.5cm}@{}}
Nama & \penulis \\
NIM & \nim \\
Departemen/Program Studi & \departemen/\programstudi \\
Fakultas & \fakultas \\
\end{tabular}

\vspace{\baselineskip}
\noindent Skripsi yang berjudul ``\judul'' merupakan hasil karya asli saya sendiri dan bukan merupakan plagiasi, fabrikasi, maupun falsifikasi, baik sebagian maupun seluruhnya.

\vspace{0.5\baselineskip}
\noindent Apabila di kemudian hari terbukti atau dapat dibuktikan bahwa skripsi ini mengandung unsur plagiasi, fabrikasi, atau falsifikasi, maka saya bersedia menerima sanksi sesuai dengan peraturan dan ketentuan yang berlaku.

\vspace{2\baselineskip}
\begin{flushright}
\begin{tabular}{c}
\kotatanggal \\
Yang membuat pernyataan, \\[2cm]
\penulis \\
NIM \nim
\end{tabular}
\end{flushright}

\clearpage
     % Pernyataan Keaslian Tulisan

    %--------------------------------------------------------------------%
%
% Ringkasan (Pedoman UM Lampiran 7a) — Bahasa Indonesia, spasi tunggal,
% maksimal 750 kata. Baris identitas: Nama-akhir, Nama-awal. Tahun. dst.
%
%--------------------------------------------------------------------%

\clearpage
\addcontentsline{toc}{chapter}{RINGKASAN}

\begin{center}
    \bfseries RINGKASAN
\end{center}

\begin{singlespace}

\noindent \penulis. \bulantahun. \textit{\judul}. Skripsi, \departemen, \fakultas, Universitas Negeri Malang. Pembimbing: (I) \pembimbingSatu\ifx\pembimbingDua\empty\else, (II) \pembimbingDua\fi

\vspace{0.5\baselineskip}
\noindent \textbf{Kata Kunci:} \katakunci%

\vspace{\baselineskip}

Model bahasa pra-latih berukuran besar telah menjadi pilihan utama pada tugas pemrosesan bahasa alami berbahasa Indonesia, termasuk \textit{extractive question answering}. Ukuran model tersebut membatasi penerapan pada perangkat dengan memori terbatas, sehingga diperlukan teknik kompresi yang tidak menuntut pelatihan ulang. Kuantisasi pascapelatihan \textit{weight-only} INT4 menjanjikan pengurangan penyimpanan hingga seperempat dibanding FP16, tetapi besarnya penurunan kinerja bergantung pada granularitas skala kuantisasi. Pada tugas \textit{extractive question answering} bahasa Indonesia, pengaruh granularitas tersebut belum diuji secara sistematis dengan uji signifikansi.

Penelitian ini mengevaluasi pengaruh granularitas kuantisasi \textit{weight-only} INT4 terhadap kinerja dan efisiensi model IndoBERT-base fase 2 pada \textit{dataset} SQuAD-id. Data disaring menjadi 70.051 baris latih, 5.997 baris validasi, dan 5.244 baris uji. Model di-\textit{fine-tune} pada presisi FP32 dengan tiga nilai \textit{seed}, lalu bobotnya disalin menjadi tiga konfigurasi, yaitu \textit{baseline} FP16, kuantisasi per-tensor, dan kuantisasi \textit{per-channel}. Kinerja diukur dengan \textit{Exact Match} dan F1, sedangkan efisiensi diukur melalui ukuran berkas model dan latensi inferensi. Perbedaan kinerja diuji dengan uji McNemar untuk EM dan \textit{paired bootstrap} 10.000 sampel ulang untuk F1, disertai koreksi Holm pada tiga perbandingan berpasangan.

Hasil penelitian menunjukkan [EM dan F1 rata-rata ketiga konfigurasi beserta selisihnya terhadap \textit{baseline}, ukuran berkas aktual dan rasionya terhadap FP16, serta nilai $p$ setelah koreksi Holm]. Temuan ini memberi gambaran mengenai pertukaran antara penghematan penyimpanan dan penurunan kinerja, sehingga dapat menjadi dasar pemilihan granularitas kuantisasi untuk penerapan IndoBERT pada perangkat dengan sumber daya terbatas. Kontribusi penelitian ini adalah bukti empiris mengenai pengaruh granularitas kuantisasi \textit{weight-only} pada tugas \textit{extractive question answering} bahasa Indonesia, beserta prosedur pengujian yang memperhitungkan variasi \textit{seed} dan koreksi perbandingan berganda.

\end{singlespace}
\clearpage
   % Ringkasan
    %--------------------------------------------------------------------%
%
% Summary (Pedoman UM Lampiran 7 — versi Bahasa Inggris dari Ringkasan)
% Spasi tunggal, maksimal 750 kata.
%
%--------------------------------------------------------------------%

\clearpage
\addcontentsline{toc}{chapter}{SUMMARY}

\begin{center}
    \bfseries SUMMARY
\end{center}

\begin{singlespace}

\noindent \penulis. \bulantahun. \textit{Weight-Only Post-Training Quantization for Indonesian Extractive Question Answering}. Thesis, \departemen, \fakultas, State University of Malang. Advisors: (I) \pembimbingSatu\ifx\pembimbingDua\empty\else, (II) \pembimbingDua\fi

\vspace{0.5\baselineskip}
\noindent \textbf{Keywords:} \katakunciEn%

\vspace{\baselineskip}

Large pre-trained language models have become the primary choice for Indonesian natural language processing tasks, including extractive question answering. Their size limits deployment on memory-constrained devices, so a compression technique that does not require retraining is needed. Weight-only INT4 post-training quantization promises up to fourfold storage reduction relative to FP16, but the resulting performance loss depends on the granularity of the quantization scale. For Indonesian extractive question answering, the effect of that granularity has not been tested systematically with significance testing.

This study evaluates the effect of INT4 weight-only quantization granularity on the performance and efficiency of the IndoBERT-base phase 2 model on the SQuAD-id dataset. The data were filtered into 70,051 training rows, 5,997 validation rows, and 5,244 test rows. The model was fine-tuned in FP32 precision with three seeds, and its weights were then copied into three configurations, namely the FP16 baseline, per-tensor quantization, and per-channel quantization. Performance was measured with Exact Match and F1, while efficiency was measured through model file size and inference latency. Performance differences were tested with the McNemar test for EM and a paired bootstrap of 10,000 resamples for F1, with Holm correction applied to three paired comparisons.

The results show [mean EM and F1 of the three configurations together with their differences from the baseline, the actual file sizes and their ratios to FP16, and the Holm-corrected $p$ values]. These findings describe the trade-off between storage savings and performance loss, and can serve as a basis for choosing quantization granularity when deploying IndoBERT on resource-constrained devices. The contribution of this study is empirical evidence on the effect of weight-only quantization granularity for Indonesian extractive question answering, together with an evaluation procedure that accounts for variation across seeds and for multiple-comparison correction.

\end{singlespace}
\clearpage
   % Summary
    %--------------------------------------------------------------------%
%
% Kata Pengantar (Pedoman UM Lampiran 7)
%
% DRAF: isi bagian dalam kurung siku dengan data sebenarnya, lalu hapus
% tanda kurung siku. Urutan ucapan terima kasih mengikuti pedoman, yaitu
% pembimbing dan program studi lebih dahulu, keluarga kemudian.
%
%--------------------------------------------------------------------%

\chapter*{Kata Pengantar}
\phantomsection
\addcontentsline{toc}{chapter}{KATA PENGANTAR}

Puji syukur penulis panjatkan kepada Tuhan Yang Maha Esa atas rahmat dan karunia-Nya. Atas kehendak-Nya penulis dapat menyelesaikan skripsi yang berjudul ``\judul''. Skripsi ini disusun untuk memenuhi salah satu persyaratan dalam menyelesaikan program \gelarprogram.

Penyusunan skripsi ini berlangsung melalui proses panjang. Pada setiap tahapnya penulis memperoleh bantuan, arahan, dan dukungan dari berbagai pihak. Karena itu, pada kesempatan ini penulis menyampaikan terima kasih kepada pihak-pihak yang berperan dalam penyelesaian skripsi ini.

\begin{enumerate}
    \item \pembimbingSatu, selaku dosen pembimbing, atas arahan, ketelitian, dan waktu yang diberikan. Bimbingan beliau membantu penulis merumuskan masalah penelitian, menyusun metode, dan menafsirkan hasil pengujian.
    \item \dekan, selaku Dekan \fakultas, atas kebijakan dan fasilitas yang mendukung pelaksanaan penelitian ini.
    \item \ketuaDepartemen, selaku Ketua Departemen \departemen, atas dukungan administratif selama masa studi penulis.
    \item Bapak Dr. Eng. Ir. Didik Dwi Prasetya, S.T., M.T., selaku Ketua \programstudi\ Universitas Negeri Malang, atas arahan mengenai ketentuan penulisan skripsi.
    \item \pembimbingSatu, selaku Dosen Pembimbing Akademik, atas nasihat dan motivasi selama masa studi.
\end{enumerate}

Penulis menyadari bahwa skripsi ini masih memiliki kekurangan. Karena itu, kritik dan saran yang membangun penulis harapkan untuk perbaikan pada penelitian selanjutnya. Penulis berharap skripsi ini bermanfaat bagi pengembangan pemrosesan bahasa alami berbahasa Indonesia.

\vspace{2\baselineskip}
\begin{flushright}
\begin{tabular}{c}
\kotatanggal \\
Penulis
\end{tabular}
\end{flushright}
     % Kata Pengantar

    \tableofcontents
    \listoftables
    \listoffigures
    \listofrumus                   % Daftar Rumus

    %----------------------------------------------------------------%
    % Bagian inti — nomor halaman arabic; kanan atas, awal bab tengah bawah.
    %----------------------------------------------------------------%
    \clearpage                     % halaman terakhir bagian awal tetap romawi
    \pagestyle{fancy}
    \pagenumbering{arabic}
    \setcounter{page}{1}

    %--------------------------------------------------------------------%
%
% Bab I — Pendahuluan
%
% @author Petra Novandi
%
%--------------------------------------------------------------------%
%
% Berisi latar belakang, rumusan masalah, tujuan penelitian, dan
% batasan masalah.
%
%--------------------------------------------------------------------%

\chapter{Pendahuluan}

\section{Latar Belakang}

Pemanfaatan model bahasa pada pengolahan bahasa alami (\textit{Natural Language Processing}) mencakup tugas klasifikasi, penerjemahan, dan \textit{question answering}, yaitu tugas menjawab pertanyaan berdasarkan informasi yang tersedia \parencite{jurafsky2025speech}. Pada \textit{extractive question answering}, sistem tidak menyusun teks baru, melainkan memilih satu rentang teks (\textit{span}) dari konteks \parencite{jurafsky2025speech}. Rentang tersebut berasal langsung dari teks sumber, sehingga model harus memahami hubungan antara pertanyaan dan konteks.

BERT (\textit{Bidirectional Encoder Representations from Transformers}) adalah model bahasa pralatih berarsitektur \textit{Transformer} yang disesuaikan pada tugas tertentu melalui \textit{fine-tuning}, termasuk \textit{question answering} \parencite{devlin2019bert}. Untuk bahasa Indonesia, IndoBERT dilatih pada korpus Indo4B dan disediakan sebagai bagian dari \textit{benchmark} IndoNLU \parencite{wilie2020indonlu}. Selain model, tugas ini memerlukan \textit{dataset} evaluasi. SQuAD (\textit{Stanford Question Answering Dataset}) dipakai pada empat penelitian kuantisasi model \textit{encoder} \parencite{bai2022towards, guo2023olive, wei2022outlier, zadeh2022mokey}. Versi terjemahannya, yaitu SQuAD-id, dipakai pada dua studi \textit{question answering} berbahasa Indonesia \parencite{ramaputra2026evaluating, suhendra2026matrik}.

Kebutuhan penyimpanan menjadi salah satu kendala penerapan model bahasa karena ukuran model bertambah sebanding dengan jumlah parameter dan jumlah bit untuk setiap parameter \parencite{gholami2021survey}. Model IndoBERT-base untuk tugas \textit{extractive question answering} memiliki 123.852.290 parameter, sehingga ukuran teoretisnya sekitar 495,4 MB pada FP32 (32 bit) dan 247,7 MB pada FP16 (16 bit). Pada sistem yang menyimpan satu model untuk setiap tugas, kebutuhan ini bertambah sebanyak model yang disediakan, misalnya sekitar 2,5 GB untuk sepuluh model pada FP16. Kondisi yang diharapkan adalah ukuran penyimpanan menurun tanpa mengurangi kemampuan model menjawab pertanyaan, tetapi dampak penurunan presisi bobot terhadap kemampuan tersebut berbeda-beda antarpenelitian.

Pada model \textit{encoder} berbahasa Indonesia, \textcite{abdurrahman2023implementing} menerapkan kuantisasi pascapelatihan (\textit{post-training quantization}, PTQ) dan \textit{quantization-aware training} (QAT) pada IndoBERT dengan bobot 16, 8, 6, dan 4 bit untuk tujuh tugas IndoNLU. F1 rata-rata pada PTQ 4-bit turun dari 77,35 pada FP32 menjadi 72,78, sedangkan QAT 4-bit mencapai 74,94, dengan ukuran model 59,75 MB dibanding 474,72 MB. \textcite{suhendra2026matrik} membandingkan model BERT berbahasa Indonesia pada SQuAD-id tanpa kuantisasi dan memakai \textit{bootstrap}, dengan F1 63,59 untuk \texttt{indobert-base-p1} dan 71,58 untuk \texttt{indolem/indobert-base-uncased}.

Pada \textit{question answering}, Ramaputra dan Joddy (\citeyear{ramaputra2026evaluating}) mengukur dampak PTQ \textit{weight-only} dengan metode AWQ dan GPTQ pada SQuAD-id menggunakan model bahasa besar \textit{decoder-only}. \textit{Exact Match} (EM) pada Qwen2.5-14B 4-bit turun dari 48,61 menjadi 44,73 (AWQ) dan 43,80 (GPTQ), sedangkan pada Qwen2.5-7B 4-bit EM berubah dari 42,05 menjadi 42,15 (AWQ) dan 42,34 (GPTQ). Pada model \textit{encoder} berbahasa Inggris, \textcite{zadeh2022mokey} memetakan bobot BERT-Large ke indeks 4-bit yang merujuk ke nilai pusat (\textit{centroid}), dan F1 pada SQuAD v1.1 berubah dari 93,15 menjadi 93,17. \textcite{bondarenko2021understanding} melaporkan rata-rata skor GLUE pada BERT-base sebesar 83,06 (FP32), 83,23 bila hanya bobot dikuantisasi 8 bit (W8A32), dan 71,03 bila bobot dan aktivasi dikuantisasi 8 bit (W8A8).

Studi pada IndoBERT melaporkan F1 rata-rata tujuh tugas dan tidak menyebut granularitas yang dipakai, sehingga pengaruh kuantisasi pada tugas \textit{span extraction} tidak dapat dipisahkan \parencite{abdurrahman2023implementing}. Kuantisasi pada SQuAD-id baru dilaporkan untuk model \textit{decoder-only} \parencite{ramaputra2026evaluating}, sedangkan studi SQuAD-id pada model \textit{encoder} tidak menerapkan kuantisasi \parencite{suhendra2026matrik}, sehingga belum diketahui apakah selisih EM tersebut berlaku pada model \textit{encoder} berparameter sekitar 124 juta. Studi yang mengkaji pengelompokan skala dilakukan pada BERT-base berbahasa Inggris dengan tugas GLUE \parencite{bondarenko2021understanding}. Granularitas relevan pada INT4 karena hanya tersedia 16 tingkat nilai, dan satu skala \textit{per-tensor} ditentukan oleh nilai absolut terbesar seluruh \textit{tensor} sehingga satu nilai ekstrem dapat memperlebar langkah kuantisasi seluruh bobot dalam \textit{tensor} itu, sedangkan skala \textit{per-channel} membatasi pengaruh tersebut pada satu baris dengan menambah jumlah skala dari 75 menjadi 133.458 \parencite{gholami2021survey}, sebagaimana dirinci pada Subbab 3.2.6. Dalam penelusuran penulis, belum ditemukan evaluasi \textit{weight-only} PTQ INT4 yang membandingkan granularitas \textit{per-tensor} dan \textit{per-channel} pada model \textit{encoder} berbahasa Indonesia untuk \textit{extractive question answering}, maupun penggunaan uji signifikansi dengan koreksi perbandingan berganda untuk selisih EM dan F1 pada evaluasi kuantisasi.

Penelitian ini menerapkan PTQ, yaitu kuantisasi yang dijalankan setelah \textit{fine-tuning} selesai sehingga tidak memerlukan pelatihan ulang \parencite{gholami2021survey}. Kuantisasi dibatasi pada bobot (\textit{weight-only}), sedangkan aktivasi tetap pada presisi asal, sehingga perubahan yang diukur hanya berasal dari representasi bobot. Presisi yang dipakai adalah INT4 (bilangan bulat 4 bit dengan 16 tingkat nilai), yang secara teoretis memperkecil ukuran bobot menjadi seperempat dari FP16, atau sekitar separuhnya bila INT4 disimpan dalam wadah 8 bit seperti pada penelitian ini. Pada INT8, \textcite{bondarenko2021understanding} melaporkan skor GLUE W8A32 sebesar 83,23 dibanding 83,06 pada FP32, sehingga penelitian ini memilih INT4 yang menurunkan presisi lebih jauh agar pengaruh granularitas lebih mungkin teramati. Pilihan ini merupakan keputusan penulis, dan INT8 tidak diuji.

Kuantisasi dijalankan dengan dua granularitas, yaitu \textit{per-tensor} (satu skala untuk seluruh \textit{tensor}, disebut \textit{coarse-grained}) dan \textit{per-channel} (satu skala untuk setiap baris keluaran, disebut \textit{fine-grained}), yang merupakan dua dari tiga tingkat granularitas pada Subbab 2.5.6. Granularitas \textit{per-group} tidak diuji karena di luar batasan penelitian. Kedua konfigurasi diterapkan pada bobot hasil \textit{fine-tuning} yang sama, sehingga selisih kinerja antarkonfigurasi hanya berasal dari perlakuan kuantisasi.

Berdasarkan uraian tersebut, penelitian ini mengevaluasi \textit{weight-only post-training quantization} INT4 pada model IndoBERT untuk tugas \textit{extractive question answering} berbahasa Indonesia menggunakan \textit{dataset} SQuAD-id. Kinerja diukur dengan EM dan F1, efisiensi diukur dengan ukuran model dan latensi inferensi, dan model pembanding adalah model FP16. Perbedaan EM dan F1 antara konfigurasi \textit{per-tensor}, \textit{per-channel}, dan \textit{baseline} diuji signifikansinya secara statistik. Hasilnya diharapkan memperlihatkan besar selisih EM dan F1 antara kedua granularitas dibandingkan jumlah skala yang disimpan, yaitu 75 untuk \textit{per-tensor} dan 133.458 untuk \textit{per-channel}. Penelitian ini berjudul \textit{Weight-Only Post-Training Quantization untuk Extractive Question Answering Bahasa Indonesia}.

\section{Rumusan Masalah}

Berdasarkan uraian pada latar belakang, penelitian ini diarahkan pada tiga pertanyaan utama, yaitu prosedur penerapan \textit{weight-only post-training quantization}, besarnya perubahan kinerja dan efisiensi yang dihasilkan, serta signifikansi statistik dari perbedaan tersebut. Pertanyaan pertama dijawab melalui prosedur yang dirinci pada Bab III, sedangkan pertanyaan kedua dan ketiga dijawab melalui eksperimen dengan metrik dan uji statistik. Dengan demikian, penelitian ini merumuskan tiga permasalahan.
\begin{enumerate}
    \item Bagaimana menerapkan \textit{weight-only post-training quantization} INT4 dengan granularitas \textit{per-tensor} dan \textit{per-channel} pada model IndoBERT untuk tugas \textit{extractive question answering} bahasa Indonesia?
    \item Seberapa besar perubahan EM, F1, ukuran model, dan latensi inferensi pada model terkuantisasi \textit{per-tensor} dan \textit{per-channel} dibandingkan model FP16 pada \textit{dataset} SQuAD-id?
    \item Apakah perbedaan EM dan F1 antarkonfigurasi signifikan secara statistik?
\end{enumerate}

\section{Tujuan Penelitian}

Berdasarkan rumusan masalah tersebut, penelitian ini memiliki tiga tujuan.
\begin{enumerate}
    \item Menerapkan \textit{weight-only post-training quantization} INT4 dengan granularitas \textit{per-tensor} dan \textit{per-channel} pada model IndoBERT untuk tugas \textit{extractive question answering} bahasa Indonesia.
    \item Mengukur perubahan EM, F1, ukuran model, dan latensi inferensi pada model terkuantisasi dibandingkan model FP16 pada \textit{dataset} SQuAD-id.
    \item Menguji signifikansi statistik perbedaan EM dan F1 antarkonfigurasi.
\end{enumerate}

\section{Manfaat Penelitian}

Manfaat penelitian terbagi menjadi manfaat teoretis (Subbab 1.4.1) dan manfaat praktis (Subbab 1.4.2).

\subsection{Manfaat Teoretis}

Secara teoretis, penelitian ini diharapkan:
\begin{enumerate}
    \item Memperkaya kajian \textit{weight-only post-training quantization} INT4 pada model \textit{encoder} berbahasa Indonesia untuk tugas \textit{extractive question answering}. Penelitian sebelumnya melaporkan IndoBERT pada tujuh tugas IndoNLU \parencite{abdurrahman2023implementing}. Kuantisasi pada \textit{question answering} dengan SQuAD-id baru dilaporkan untuk model bahasa besar \textit{decoder-only} \parencite{ramaputra2026evaluating}.
    \item Menyediakan data pembanding antara granularitas \textit{per-tensor} dan \textit{per-channel} pada model, \textit{dataset}, dan presisi yang sama.
    \item Memberikan data empiris mengenai pertukaran antara kinerja (EM dan F1) dan ukuran model pada kuantisasi INT4, termasuk selisih antara ukuran teoretis (hasil hitungan dari jumlah bit per parameter) dan ukuran berkas (ukuran riil pada penyimpanan).
    \item Menunjukkan konsistensi arah selisih kinerja antarkonfigurasi pada tiga nilai \textit{seed}, disertai uji signifikansi berpasangan dan koreksi perbandingan berganda.
\end{enumerate}

\subsection{Manfaat Praktis}

Secara praktis, penelitian ini diharapkan:
\begin{enumerate}
    \item Menjadi bahan pertimbangan bagi pengembang sistem \textit{question answering} berbahasa Indonesia dalam memilih konfigurasi kuantisasi bobot IndoBERT berdasarkan perubahan EM, F1, dan ukuran model.
    \item Menjadi bahan pembanding bagi peneliti selanjutnya untuk pengembangan skema kuantisasi lain, misalnya granularitas \textit{per-group}, kuantisasi bobot dan aktivasi sekaligus, atau \textit{quantization-aware training}.
    \item Memberikan gambaran kuantitatif penghematan ruang penyimpanan bagi institusi yang mempertimbangkan penerapan model bahasa pada infrastruktur dengan ruang penyimpanan terbatas.
\end{enumerate}

\section{Batasan Masalah}

Batasan masalah dalam penelitian ini terdiri atas sepuluh butir.
\begin{enumerate}
    \item Model yang digunakan adalah IndoBERT-base fase 2 \parencite{wilie2020indonlu}, yaitu model bahasa pralatih umum yang tidak dikhususkan untuk tugas \textit{question answering}. Penyesuaian ke tugas \textit{extractive question answering} dilakukan melalui \textit{fine-tuning} pada penelitian ini, dan tidak ada model dasar lain yang dibandingkan.
    \item \textit{Dataset} yang digunakan adalah SQuAD-id \parencite{muis2020sequence}, yaitu terjemahan mesin SQuAD v2.0 ke bahasa Indonesia, dalam versi yang didistribusikan SEACrowd \parencite{lovenia2024seacrowd}. Setiap baris berisi satu pasangan pertanyaan dan konteks. Kolom yang dipakai adalah konteks, pertanyaan, jawaban (teks dan posisi awal hasil perbaikan), dan judul artikel untuk pembagian data. Data validasi dibentuk dari sebagian data latih.
    \item Sampel yang tidak diikutsertakan terdiri atas dua jenis, yaitu pertanyaan tanpa jawaban pada konteks (bagian dari rancangan SQuAD v2.0) dan pertanyaan terjawab yang jawabannya tidak dapat dilacak pada konteks (galat \textit{dataset}). Metrik hanya dihitung pada sampel yang memiliki jawaban, sehingga penelitian ini tidak mengukur kemampuan model memutuskan bahwa suatu pertanyaan tidak dapat dijawab. Jumlah setiap jenis disajikan pada Subbab 3.2.4.
    \item Mutu terjemahan mesin tidak dinilai melalui evaluasi manusia (\textit{human evaluation}) pada penelitian ini. Oleh karena itu, hasil penelitian tidak dapat digeneralisasi ke \textit{dataset} berbahasa Indonesia yang ditulis langsung oleh penutur jati.
    \item Kuantisasi yang diterapkan adalah \textit{weight-only post-training quantization} pada bobot model hasil \textit{fine-tuning}, tanpa pelatihan ulang dan tanpa kuantisasi aktivasi. Kuantisasi bersifat tersimulasi, yaitu bobot dikembalikan ke presisi \textit{floating point} sebelum komputasi, sehingga penelitian ini tidak mengklaim percepatan inferensi. Model pembanding adalah model FP16.
    \item Kuantisasi diterapkan pada matriks bobot lapisan linear dan \textit{embedding}. Kepala jawaban, \textit{bias}, dan normalisasi lapisan tetap pada presisi asal.
    \item Presisi numerik dibatasi pada INT4, sehingga perbandingan antargranularitas dilakukan pada jumlah bit yang sama.
    \item Granularitas dibatasi pada \textit{per-tensor} (satu skala untuk setiap \textit{tensor}) dan \textit{per-channel} (satu skala untuk setiap baris keluaran). Jumlah \textit{tensor} dan skala disajikan pada Subbab 3.2.6.
    \item Kinerja diukur dengan EM dan F1, sedangkan efisiensi diukur dengan ukuran model dan latensi inferensi. Hanya EM dan F1 yang diuji signifikansinya secara statistik, sedangkan latensi dilaporkan secara deskriptif.
    \item Eksperimen dijalankan pada satu jenis \textit{graphics processing unit} (GPU), yaitu NVIDIA Tesla T4, dengan tiga nilai \textit{seed}. Penelitian ini tidak menerapkan \textit{quantization-aware training}, tidak mengembangkan \textit{kernel} khusus untuk presisi rendah, dan tidak menguji model pada perangkat dengan sumber daya terbatas.
\end{enumerate}

    %--------------------------------------------------------------------%
%
% Bab II — Kajian Pustaka
%
% @author Petra Novandi, disesuaikan untuk penelitian kuantisasi
%
%--------------------------------------------------------------------%

\chapter{Kajian Pustaka}

\section{Tugas \textit{Extractive Question Answering}}

\textit{Question answering} adalah tugas pengolahan bahasa alami yang meminta sistem menjawab pertanyaan berdasarkan informasi yang tersedia \parencite{jurafsky2025speech}. Tugas ini termasuk dalam \textit{machine reading comprehension}, yaitu menjawab pertanyaan berdasarkan teks tertentu, dan dapat dipandang sebagai masalah pembelajaran terbimbing dengan masukan berupa konteks dan pertanyaan dalam bahasa alami \parencite{wang2022extractive}. Pada pendekatan \textit{extractive} (disebut juga \textit{span prediction}), jawaban tidak dihasilkan dari awal, melainkan dipilih sebagai rentang teks (\textit{span}) di dalam konteks \parencite{jurafsky2025speech, wang2022extractive}. Penelitian ini memakai bentuk \textit{single-span}, yaitu hanya satu rentang yang dipilih sebagai jawaban.

Pada bentuk \textit{single-span}, model memilih satu urutan \textit{token} yang berurutan dari konteks sebagai jawaban atas pertanyaan \parencite{wang2022extractive}. Masukan tugas ini berupa pasangan pertanyaan dan konteks, sedangkan keluarannya berupa posisi \textit{token} awal dan posisi \textit{token} akhir \textit{span} jawaban \parencite{devlin2019bert}. Dengan formulasi tersebut, \textit{extractive question answering} dapat dipandang sebagai dua klasifikasi posisi pada urutan \textit{token} konteks, sehingga kualitas jawaban dapat dibandingkan langsung dengan jawaban acuan melalui kecocokan teks dan \textit{token overlap}. Perbandingan ini memakai \textit{Exact Match} (EM), yaitu proporsi jawaban yang persis sama dengan jawaban acuan, dan F1, yaitu ukuran \textit{token overlap} antara jawaban prediksi dan jawaban acuan \parencite{rajpurkar2016squad}.

Model \textit{extractive question answering} tersusun dari empat komponen, yaitu \textit{embedding}, ekstraksi fitur, interaksi antara konteks dan pertanyaan, serta prediksi jawaban \parencite{wang2022extractive}. Ekstraksi fitur memakai jaringan saraf berulang (\textit{recurrent neural network}), jaringan saraf konvolusional (\textit{convolutional neural network}), atau struktur \textit{Transformer} berbasis \textit{multi-head self-attention}, sedangkan model bahasa pralatih seperti BERT dipasangkan dengan \textit{pointer head} untuk menentukan rentang jawaban \parencite{wang2022extractive}. Pada BERT, \textit{fine-tuning} ke tugas ini hanya menambahkan satu lapisan keluaran \parencite{raikwar2024five}. Model yang mampu memodelkan hubungan antara pertanyaan dan konteks diperlukan untuk menentukan kedua posisi tersebut, sehingga pembahasan dilanjutkan pada arsitektur \textit{Transformer}.

SQuAD (\textit{Stanford Question Answering Dataset}) diperkenalkan pada 2016 dan berperan penting dalam pelatihan dan evaluasi model \textit{machine comprehension} \parencite{rajpurkar2016squad, raikwar2024five}. Versi 2.0 memuat pertanyaan yang tidak memiliki jawaban pada konteks \parencite{rajpurkar2018squad}, dan pertanyaan semacam itu juga dicatat sebagai salah satu masalah pada \textit{dataset} \textit{single-span} \parencite{wang2022extractive}. Penelitian ini hanya memakai pertanyaan yang memiliki jawaban, karena kemampuan model memutuskan bahwa suatu pertanyaan tidak dapat dijawab tidak termasuk yang diuji (Subbab 1.5).

Pada bahasa Inggris, tugas ini telah dikaji dengan membandingkan lima model berbasis BERT yang di-\textit{fine-tune} pada SQuAD \parencite{raikwar2024five}. Pada bahasa Indonesia, \textit{extractive question answering} dievaluasi pada SQuAD-id \parencite{muis2020sequence, ramaputra2026evaluating, suhendra2026matrik} dan pada FacQA, yaitu tugas \textit{span extraction} berita di dalam \textit{benchmark} IndoNLU \parencite{wilie2020indonlu}. Penelitian ini memakai SQuAD-id, yang dijelaskan pada Subbab 2.3.

\section{Model Bahasa Pralatih BERT dan IndoBERT}

\subsection{Arsitektur dan Cara Kerja BERT}

\textit{Transformer} adalah arsitektur yang memproses urutan \textit{token} dengan mekanisme \textit{attention} tanpa rekurensi \parencite{vaswani2017attention}. Setiap lapisan memakai matriks proyeksi \textit{query}, \textit{key}, \textit{value}, dan keluaran, diikuti lapisan \textit{feed-forward}. Matriks-matriks inilah, bersama \textit{embedding}, yang menjadi bobot yang dikuantisasi pada penelitian ini.

BERT (\textit{Bidirectional Encoder Representations from Transformers}) adalah model bahasa yang dibangun dari tumpukan \textit{encoder} \textit{Transformer} dan dipralatih pada korpus teks tanpa label melalui dua tugas, yaitu \textit{masked language modeling} (MLM) dan \textit{next sentence prediction} (NSP) \parencite{devlin2019bert}. Pada MLM, 15\% \textit{token} dipilih secara acak sebagai target prediksi dan, pada 80\% di antaranya, diganti dengan \textit{token} khusus \textit{[MASK]}, lalu model memprediksi \textit{token} target dengan memanfaatkan konteks di kiri dan kanan sekaligus, sehingga representasinya bersifat dua arah (\textit{bidirectional}). Pada NSP, model menerima sepasang kalimat dan menentukan apakah kalimat kedua merupakan kelanjutan dari kalimat pertama atau diambil secara acak dari korpus. Kombinasi kedua tugas ini membuat BERT memperoleh pengetahuan bahasa umum yang kemudian diadaptasi ke tugas spesifik melalui \textit{fine-tuning} \parencite{devlin2019bert}.

Teks masukan terlebih dahulu ditokenisasi dengan \textit{WordPiece}, yaitu metode yang memecah kata jarang menjadi potongan subkata (\textit{subword}) \parencite{devlin2019bert}. Tokenisasi menambahkan \textit{token} khusus \textit{[CLS]} di awal urutan dan \textit{[SEP]} sebagai pemisah antarsegmen, serta menghasilkan tiga keluaran, yaitu \textit{input\_ids} (representasi numerik \textit{token}), \textit{attention\_mask} (penanda \textit{token} asli terhadap \textit{padding}), dan \textit{token\_type\_ids} (penanda segmen). Setiap \textit{token} dipetakan ke tiga \textit{embedding}, yaitu \textit{embedding} \textit{token}, segmen, dan posisi, yang dijumlahkan sebelum masuk ke lapisan \textit{encoder} \parencite{devlin2019bert}. Pada setiap lapisan \textit{encoder}, mekanisme \textit{self-attention} memungkinkan sebuah \textit{token} berinteraksi dengan \textit{token} lain dalam urutan yang sama, termasuk yang posisinya berjauhan, sehingga setiap \textit{token} memperoleh representasi kontekstual \parencite{vaswani2017attention}.

Penyesuaian ke tugas \textit{extractive question answering} dilakukan melalui \textit{fine-tuning}. Pertanyaan dan konteks digabung menjadi satu urutan masukan yang dipisahkan oleh \textit{[SEP]}, lalu satu lapisan linear menghasilkan skor posisi awal dan posisi akhir untuk setiap \textit{token}, dan \textit{span} dengan jumlah skor awal dan akhir tertinggi dipilih sebagai jawaban \parencite{devlin2019bert}. Lapisan tambahan ini hanya berukuran $768 \times 2$, yaitu 1.536 bobot, sehingga ukuran model hasil \textit{fine-tuning} ditentukan oleh \textit{encoder} dan \textit{embedding}, dan ukuran itu menentukan kebutuhan penyimpanan. Karena itu, kompresi model dibahas pada Subbab 2.4.

\subsection{Model IndoBERT-base yang Digunakan}

IndoBERT adalah model bahasa pralatih berbahasa Indonesia berarsitektur BERT yang disediakan sebagai bagian dari \textit{benchmark} IndoNLU dan dipralatih pada korpus Indo4B \parencite{wilie2020indonlu}. Model yang digunakan pada penelitian ini adalah IndoBERT-base fase 2. Konfigurasinya terdiri atas 12 lapisan \textit{encoder}, ukuran \textit{hidden} 768, 12 \textit{head} perhatian, ukuran lapisan \textit{feed-forward} 3.072, kosakata 50.000 \textit{token}, dan panjang urutan maksimum 512 \textit{token}. Model untuk tugas \textit{question answering} tidak memakai lapisan \textit{pooler}, sehingga jumlah parameternya 123.852.290 berdasarkan penghitungan pada model yang digunakan. Dengan lapisan \textit{pooler}, jumlahnya menjadi 124.442.882.

Nama model yang dipakai dinyatakan secara eksplisit pada Subbab 3.2.5 karena kinerja pada SQuAD-id berbeda antarmodel berbahasa Indonesia. Pada SQuAD-id, \textcite{suhendra2026matrik} melaporkan F1 sebesar 63,59 untuk model \texttt{indobert-base-p1}. Untuk model \texttt{indolem/indobert-base-uncased}, F1 yang dilaporkan adalah 71,58. IndoBERT-base dipralatih dalam dua fase. Fase 1 memakai panjang urutan maksimum 128 \textit{token}, dan fase 2 melanjutkan pralatih dengan panjang urutan maksimum 512 \textit{token} \parencite{wilie2020indonlu}. Penelitian ini memakai fase 2 (\texttt{indobert-base-p2}) karena konteks SQuAD-id diproses dalam jendela hingga 384 \textit{token} (Bab III), lebih panjang daripada 128 \textit{token} pada fase 1. Pada FacQA, skor IndoBERT-base naik dari 53,45 pada fase 1 menjadi 57,06 pada fase 2 \parencite{wilie2020indonlu}.

\section{Karakteristik \textit{Dataset} SQuAD-id}

SQuAD-id adalah \textit{dataset} \textit{extractive question answering} berbahasa Indonesia yang dibentuk dari terjemahan mesin SQuAD v2.0 ke bahasa Indonesia oleh \textcite{muis2020sequence} untuk penelitian pembangkit pertanyaan otomatis. Setiap baris memuat satu pasangan pertanyaan dan konteks dengan kolom konteks, pertanyaan, jawaban (teks dan posisi awal), judul artikel, dan penanda pertanyaan tanpa jawaban. \textit{Dataset} asalnya, SQuAD, memuat jawaban berupa \textit{span} teks dari paragraf Wikipedia \parencite{rajpurkar2016squad}.

Berkas yang dipakai hanya menyediakan data latih dan data uji tanpa data validasi, sehingga data validasi dibentuk dari data latih pada tahap praproses (Subbab 3.2.3). Jumlah data dan panjang teks dirangkum pada Tabel~\ref{tab:karakteristik-dataset}. Tabel tersebut dipakai untuk menunjukkan jumlah baris yang tidak dapat dipakai, baik karena pertanyaannya tidak memiliki jawaban maupun karena teks jawabannya tidak dapat dilacak pada konteks.

\begin{table}[htbp]
  \centering
  \small
  \renewcommand{\arraystretch}{1.2}
  \caption{Karakteristik \textit{dataset} SQuAD-id}
  \label{tab:karakteristik-dataset}
  \begin{tabularx}{\textwidth}{@{}Xrr@{}}
    \hline
    \textbf{Karakteristik} & \textbf{Data latih} & \textbf{Data uji} \\
    \hline
    Jumlah baris & 130.319 & 11.873 \\
    Artikel unik & 441 & 35 \\
    Konteks unik & 19.029 & 1.204 \\
    Pertanyaan tanpa jawaban & 43.498 (33,38\%) & 5.945 (50,07\%) \\
    Pertanyaan terjawab & 86.821 & 5.928 \\
    Jawaban tidak terlacak pada konteks setelah perbaikan & 16.705 & 676 \\
    Duplikat konteks dan pertanyaan & 65 & 8 \\
    Baris yang dipakai pada penelitian ini & 70.051 (53,75\%) & 5.244 (44,17\%) \\
    Median panjang konteks (kata) & 98 & 101 \\
    Median panjang pertanyaan (kata) & 8 & 8 \\
    \hline
  \end{tabularx}
\end{table}

Dua batasan \textit{dataset} perlu dicatat. Pertama, SQuAD-id berasal dari terjemahan mesin, sehingga dapat membawa pengaruh proses terjemahan pada kelancaran kalimat dan ketepatan istilah. Kedua, hasil pada satu \textit{dataset} hasil terjemahan tidak dapat langsung digeneralisasi ke domain atau sumber teks lain.

\section{Teknik Kompresi Model}

Kompresi model mencakup tiga kelompok teknik untuk mengurangi kebutuhan sumber daya, yaitu \textit{pruning}, \textit{knowledge distillation}, dan kuantisasi \parencite{gholami2021survey, han2016deep, hinton2015distilling}.

\textit{Pruning} menghapus bobot atau koneksi yang kontribusinya kecil terhadap keluaran model \parencite{han2016deep}. Struktur bobot yang dihasilkan tidak beraturan, sehingga penghematan waktu komputasi bergantung pada dukungan perangkat keras terhadap operasi \textit{sparse}, dan model perlu dilatih ulang untuk memulihkan akurasi \parencite{han2016deep}.

\textit{Knowledge distillation} melatih model berukuran lebih kecil (\textit{student}) agar meniru distribusi keluaran model besar (\textit{teacher}) \parencite{hinton2015distilling}. Teknik ini menghasilkan model baru dengan arsitektur berbeda, bukan versi terkompresi dari model yang sudah ada, sehingga memerlukan proses pelatihan penuh.

Kuantisasi menurunkan presisi representasi numerik bobot atau aktivasi \parencite{gholami2021survey}. Berbeda dengan dua teknik sebelumnya, kuantisasi tidak mengubah arsitektur maupun jumlah parameter dan dapat diterapkan pada model yang telah selesai dilatih tanpa pelatihan ulang \parencite{gholami2021survey}. Bobot hasil \textit{fine-tuning} tetap dapat dipertahankan sebagai acuan, dan \textit{error} yang muncul dapat dikendalikan melalui pilihan cakupan dan granularitas.

Besarnya pengurangan ukuran dapat dihitung langsung dari jumlah parameter. Model dengan 123.852.290 parameter berukuran 495,4 MB pada format \textit{floating point} 32 bit (FP32), 247,7 MB pada \textit{floating point} 16 bit (FP16), 123,9 MB pada bilangan bulat 8 bit (INT8), dan 61,9 MB pada bilangan bulat 4 bit (INT4) apabila setiap parameter menempati 4, 2, 1, dan 0,5 byte. Angka tersebut adalah ukuran teoretis, dan 1 MB pada penelitian ini sama dengan $10^{6}$ byte.

Hasil kuantisasi 4-bit pada model \textit{encoder} dan model bahasa besar dilaporkan sebagai berikut. Pada BERT-Large untuk SQuAD v1.1, kuantisasi bobot 4-bit mengubah F1 dari 93,15 menjadi 93,17 \parencite{zadeh2022mokey}. Pada IndoBERT, kuantisasi pascapelatihan 4-bit menurunkan F1 rata-rata tujuh tugas dari 77,35 menjadi 72,78 \parencite{abdurrahman2023implementing}. Pada SQuAD-id dengan model bahasa besar, EM Qwen2.5-14B 4-bit turun dari 48,61 menjadi 44,73 (AWQ) dan 43,80 (GPTQ), sedangkan EM Qwen2.5-7B 4-bit berubah dari 42,05 menjadi 42,15 (AWQ) dan 42,34 (GPTQ) \parencite{ramaputra2026evaluating}. Perubahan F1 berkisar dari naik 0,02 poin hingga turun 4,57 poin, dan perubahan EM dari naik 0,29 poin hingga turun 4,81 poin, sehingga dampak kuantisasi perlu diukur pada model dan tugas yang spesifik, termasuk \textit{extractive question answering} berbahasa Indonesia.

\section{Konsep dan Skema Kuantisasi pada Jaringan Saraf Tiruan}

Dalam jaringan saraf, kuantisasi diterapkan pada \textit{tensor} bobot dan aktivasi yang tersimpan dalam format \textit{floating point}. Proses ini beroperasi dua arah: kuantisasi memetakan nilai riil ke bilangan bulat berpresisi rendah, sedangkan \textit{dequantization} mengembalikan nilai tersebut ke bentuk riil untuk komputasi. Ruang desain kuantisasi ditunjukkan pada Gambar~\ref{fig:alur-kuantisasi} dalam tujuh tahapan, yaitu representasi numerik awal, skema pemetaan simetris atau asimetris, cakupan bobot dan aktivasi, waktu penerapan, granularitas, representasi memori \textit{bit-packing}, serta \textit{dequantization} dan rekonstruksi inferensi. Subbab 2.5.1 sampai 2.5.7 membahas tahapan tersebut, dengan \textit{dequantization} dijelaskan pada Subbab 2.5.2 dan 2.5.7.

\begin{figure}[htbp]
  \centering
  \includegraphics[width=0.90\textwidth]{"resources/alur kuantisasi.png"}
  \caption{Tahapan dan ruang desain kuantisasi pada jaringan saraf}
  \label{fig:alur-kuantisasi}
\end{figure}

\subsection{Format Representasi Numerik FP16, INT8, dan INT4}

Presisi numerik menunjukkan jumlah bit untuk menyimpan setiap nilai. Format \textit{floating point} seperti FP32 dan FP16 memisahkan bit menjadi tanda, eksponen, dan mantisa, sehingga mampu mewakili rentang nilai hingga $\pm 3{,}4 \times 10^{38}$ pada FP32 (Tabel~\ref{tab:format-numerik}). Format bilangan bulat seperti INT8 dan INT4 hanya menyimpan nilai diskret dengan rentang tetap dan memerlukan faktor skala untuk memetakan nilai riil \parencite{gholami2021survey}. INT8 menyediakan 256 tingkat nilai dan INT4 hanya 16 tingkat, sehingga pada rentang bobot yang sama, langkah kuantisasi INT4 lebih lebar daripada INT8 dan \textit{error} pembulatannya lebih besar \parencite{gholami2021survey}. Susunan bit dan rentang nilai keempat format dirangkum pada Tabel~\ref{tab:format-numerik}. Bilangan bulat bertanda disimpan dalam bentuk komplemen dua, dan pada kuantisasi simetris yang dipakai penelitian ini INT4 hanya memakai $-7$ sampai 7 (15 tingkat nilai) dengan nilai $-8$ tidak dipakai (Subbab 2.5.3). Pembagian bit FP16, INT8, dan INT4 ditunjukkan pada Gambar~\ref{fig:layout-fp16} sampai Gambar~\ref{fig:layout-int4}.

Pada format \textit{floating point}, satu nilai disusun dari tiga bagian. Bit tanda menyimpan apakah nilai positif (0) atau negatif (1). Bit eksponen menyimpan pangkat dua yang menggeser titik desimal, sehingga menentukan seberapa besar atau kecil nilai yang dapat diwakili. Bit mantisa menyimpan digit pecahan di belakang angka 1 awal, sehingga menentukan ketelitian nilai. Pada FP16 terdapat 1 bit tanda, 5 bit eksponen, dan 10 bit mantisa, dan nilai bilangan normalnya dihitung dengan $(-1)^{s} \times 2^{e-15} \times (1 + m/1024)$, dengan $s$ bit tanda, $e$ nilai eksponen, $m$ nilai mantisa, dan 15 bias eksponen.

Sebagai contoh, bobot 0,30 ditulis pada FP16 sebagai $1{,}2 \times 2^{-2}$. Eksponen $-2$ disimpan sebagai $-2 + 15 = 13$, yaitu \texttt{01101} dalam biner. Pecahan 0,2 disimpan sebagai $0{,}2 \times 1024 \approx 205$, yaitu \texttt{0011001101} dalam biner. Dengan bit tanda 0, susunan bitnya adalah \texttt{0 01101 0011001101}, yang dibaca kembali sebagai $2^{-2} \times (1 + 205/1024) \approx 0{,}30005$. Selisih 0,00005 terhadap nilai asli muncul karena mantisa hanya 10 bit.

Format bilangan bulat tidak memiliki bit eksponen dan bit mantisa. Seluruh bit menyimpan satu bilangan bulat dengan bit tertinggi sebagai tanda, sehingga jarak antar nilai yang dapat disimpan selalu sama. Karena itu, nilai riil harus dibagi dengan faktor skala terlebih dahulu agar muat pada rentang bilangan bulat (Subbab 2.5.2).

\begin{table}[htbp]
  \centering
  \small
  \renewcommand{\arraystretch}{1.2}
  \caption{Susunan bit dan rentang nilai format numerik \parencite{gholami2021survey}}
  \label{tab:format-numerik}
  \setlength{\tabcolsep}{4pt}
  \begin{tabularx}{\textwidth}{@{}lcccccc@{}}
    \hline
    \textbf{Format} & \textbf{Total bit} & \textbf{Tanda} & \textbf{Eksponen} & \textbf{Mantisa} & \textbf{Rentang nilai} & \textbf{Presisi} \\
    \hline
    FP32 & 32 & 1 & 8 & 23 & $\pm 3{,}4 \times 10^{38}$ & 7 digit desimal \\
    FP16 & 16 & 1 & 5 & 10 & $\pm 6{,}6 \times 10^{4}$ & 3 digit desimal \\
    INT8 & 8 & 1 & -- & -- & $-128$ sampai 127 & 256 tingkat \\
    INT4 & 4 & 1 & -- & -- & $-8$ sampai 7 & 16 tingkat \\
    \hline
  \end{tabularx}
\end{table}

\begin{figure}[htbp]
  \centering
  \includegraphics[width=0.86\textwidth]{"resources/IEEE FP16 Bit Layout Diagram.png"}
  \caption{Susunan bit pada format \textit{floating point} IEEE FP16}
  \label{fig:layout-fp16}
\end{figure}

\begin{figure}[htbp]
  \centering
  \includegraphics[width=0.86\textwidth]{"resources/8-Bit Signed Integer Layout.png"}
  \caption{Susunan bit pada bilangan bulat bertanda 8-bit (INT8)}
  \label{fig:layout-int8}
\end{figure}

\begin{figure}[htbp]
  \centering
  \includegraphics[width=0.86\textwidth]{"resources/4-Bit Signed Integer Layout.png"}
  \caption{Susunan bit pada bilangan bulat bertanda 4-bit (INT4)}
  \label{fig:layout-int4}
\end{figure}

Ukuran teoretis model sebanding dengan jumlah bit per parameter. Dibandingkan FP16, ukuran teoretis INT4 adalah 25\% apabila seluruh parameter disimpan pada 4 bit dan dipadatkan. Angka itu adalah batas teoretis, bukan ukuran berkas tanpa \textit{bit-packing} (Subbab 2.5.7).

\subsection{Pemetaan Nilai Riil ke Bilangan Bulat}

Kuantisasi pada jaringan saraf diperkenalkan sebagai cara menggantikan komputasi \textit{floating point} dengan aritmetika bilangan bulat pada inferensi tanpa mengubah arsitektur model \parencite{jacob2018quantization}. Pemetaan nilai riil $x$ ke bilangan bulat $x_q$ berbentuk \parencite{jacob2018quantization}:
\begin{equation}
  x_q = \mathrm{round}\left(\frac{x}{s}\right) + z
  \label{eq:kuantisasi}
\end{equation}
\noindent dengan $s$ faktor skala dan $z$ \textit{zero point} (bilangan bulat yang mewakili nilai riil nol). Hasilnya dibatasi pada rentang bilangan bulat yang tersedia. \textit{Dequantization} mengembalikan nilai riil melalui $\hat{x} = s\,(x_q - z)$. Selisih antara $x$ dan $\hat{x}$ adalah \textit{error} kuantisasi yang ditentukan oleh lebar langkah $s$. Skema simetris dan asimetris (Subbab 2.5.3), cakupan (Subbab 2.5.4), dan granularitas (Subbab 2.5.6) pada dasarnya adalah cara berbeda menentukan $s$ dan $z$.

Penerapan kuantisasi pada model terlatih dimulai dengan menentukan rentang nilai setiap \textit{tensor}. Pada kuantisasi bobot saja, rentang dihitung dari bobot itu sendiri sehingga tidak memerlukan data kalibrasi, sedangkan kuantisasi aktivasi memerlukan data untuk mengamati rentang aktivasi \parencite{gholami2021survey}. Dari rentang tersebut dihitung $s$ dan $z$, lalu nilai riil dipetakan ke bilangan bulat dengan Persamaan~\ref{eq:kuantisasi}. Nilai terkuantisasi disimpan beserta skalanya, dan saat inferensi nilai tersebut dikembalikan melalui \textit{dequantization} sebelum komputasi.

\subsection{Skema Kuantisasi Simetris dan Asimetris}

Skema simetris menetapkan \textit{zero point} sama dengan nol dan membatasi bilangan bulat $b$-bit pada $-q_{\max}$ sampai $q_{\max}$, dengan $q_{\max} = 2^{b-1}-1$ \parencite{gholami2021survey}:
\begin{equation}
  x_q = \mathrm{round}\left(\frac{x}{s}\right)
  \label{eq:simetris}
\end{equation}
\begin{equation}
  s = \frac{\alpha}{q_{\max}}
  \label{eq:skala}
\end{equation}
\noindent dengan $\alpha$ nilai absolut terbesar pada cakupan yang dipilih. Untuk $b = 4$, $q_{\max} = 7$, sehingga bilangan bulat berada pada $-7$ sampai 7 dan nilai $-8$ tidak dipakai. \textit{Dequantization} dilakukan dengan $\hat{x} = x_q \cdot s$.

Sebagai contoh, bobot 0,30 pada \textit{tensor} dengan $\alpha = 0{,}8$ disimpan pada tiga format lalu dibaca kembali. Pada FP16, bobot ditulis langsung sebagai bit tanda, eksponen, dan mantisa (Subbab 2.5.1). Pada INT8, $q_{\max} = 127$ sehingga $s = 0{,}8/127 \approx 0{,}0063$, lalu $0{,}30/0{,}0063 \approx 47{,}6$ dibulatkan menjadi 48 dan dibaca kembali sebagai $48 \times 0{,}0063 \approx 0{,}3024$. Pada INT4, $q_{\max} = 7$ sehingga $s = 0{,}8/7 \approx 0{,}1143$, lalu $0{,}30/0{,}1143 \approx 2{,}63$ dibulatkan menjadi 3 dan dibaca kembali sebagai $3 \times 0{,}1143 \approx 0{,}3429$. Hasil ketiganya dirangkum pada Tabel~\ref{tab:contoh-format}. \textit{Error} INT4 paling besar karena langkah kuantisasi $s$ lebih lebar 18,1 kali daripada INT8 ($127/7$), dan \textit{error} pembulatan setiap bobot paling besar setengah dari $s$, yaitu 0,0031 pada INT8 dan 0,0571 pada INT4.

\begin{table}[htbp]
  \centering
  \footnotesize
  \setlength{\tabcolsep}{4pt}
  \renewcommand{\arraystretch}{1.2}
  \caption{Contoh penyimpanan bobot 0,30 pada FP16, INT8, dan INT4}
  \label{tab:contoh-format}
  \begin{tabularx}{\textwidth}{@{}>{\centering\arraybackslash}p{1.5cm}>{\centering\arraybackslash}p{1.9cm}>{\centering\arraybackslash}p{3.7cm}>{\centering\arraybackslash}X>{\centering\arraybackslash}X@{}}
    \hline
    \textbf{Format} & \textbf{Faktor skala $s$} & \textbf{Nilai tersimpan} & \textbf{Nilai dibaca kembali} & \textbf{\textit{Error}} \\
    \hline
    FP16 & -- & \mbox{\texttt{0 01101 0011001101}} & 0,30005 & 0,00005 \\
    INT8 & 0,0063 & 48 & 0,3024 & 0,0024 \\
    INT4 & 0,1143 & 3 & 0,3429 & 0,0429 \\
    \hline
  \end{tabularx}
\end{table}

Pada skema asimetris, rentang dipetakan dari nilai terkecil hingga terbesar menggunakan seluruh tingkat dengan tambahan parameter $z$. Misalkan satu \textit{tensor} bernilai $-0{,}8$ sampai 0,6. Skema simetris memakai $\alpha = 0{,}8$ dengan $s \approx 0{,}114$, dan dua tingkat positif tertinggi tidak terpakai karena data hanya sampai 0,6. Skema asimetris memakai seluruh 16 tingkat dengan $s = 1{,}4/15 \approx 0{,}093$ dan $z = 1$. Perbedaan keduanya ditunjukkan pada Gambar~\ref{fig:pemetaan-simetris} dan Gambar~\ref{fig:pemetaan-asimetris}.

\begin{figure}[htbp]
  \centering
  \includegraphics[width=0.80\textwidth]{"resources/Pemetaan Rentang FP32 ke INT8.png"}
  \caption{Pemetaan rentang pada kuantisasi simetris}
  \label{fig:pemetaan-simetris}
\end{figure}

\begin{figure}[htbp]
  \centering
  \includegraphics[width=0.80\textwidth]{"resources/Asymmetric FP32 to INT8 Mapping.png"}
  \caption{Pemetaan rentang pada kuantisasi asimetris}
  \label{fig:pemetaan-asimetris}
\end{figure}

Skema simetris dipilih karena bobot model terlatih terkonsentrasi di sekitar nol, sehingga \textit{zero point} tidak diperlukan \parencite{gholami2021survey}. Aktivasi setelah fungsi aktivasi yang hanya menghasilkan nilai non-negatif tidak simetris terhadap nol, sehingga skema asimetris dipakai untuk aktivasi tersebut \parencite{gholami2021survey}.

\subsection{Cakupan Kuantisasi \textit{Weight-Only} dan \textit{Joint Weight-Activation}}

Kuantisasi dapat diterapkan pada bobot saja atau pada bobot dan aktivasi sekaligus. Notasi W$x$A$y$ menunjukkan bobot $x$-bit dan aktivasi $y$-bit, misalnya W8A8 untuk bobot dan aktivasi 8-bit \parencite{bondarenko2021understanding}. Pada BERT-base, \textcite{bondarenko2021understanding} melaporkan rata-rata skor GLUE sebesar 83,06 (FP32), 71,03 (W8A8), 70,25 (W32A8), dan 83,23 (W8A32). Penurunan skor muncul ketika aktivasi dikuantisasi, sedangkan kuantisasi bobot saja menghasilkan skor 83,23, 0,17 poin di atas FP32. Pada \textit{weight-only} PTQ, bobot dikuantisasi setelah pelatihan dan aktivasi tetap pada presisi asalnya.

Cakupan bobot yang dikuantisasi juga memengaruhi hasil. Model IndoBERT-base untuk \textit{question answering} terdiri atas tiga kelompok parameter. Parameter \textit{embedding} berjumlah 38.794.752 (31,3\%), parameter lapisan linear berjumlah 84.934.656 (68,6\%) yang meliputi proyeksi atensi dan lapisan \textit{feed-forward}, dan sisanya 122.882 parameter (0,1\%) berupa \textit{answer head}, \textit{bias}, dan normalisasi lapisan. Penelitian ini mengkuantisasi matriks bobot lapisan linear dan \textit{embedding} (Subbab 1.5). Apabila seluruh matriks bobot tersebut dikuantisasi INT4 dan disimpan dalam wadah satu \textit{byte}, ukuran bobot terkuantisasi adalah 50\% dari ukuran FP16, dan 25\% bila dipadatkan, di luar parameter yang tidak dikuantisasi dan faktor skala (Subbab 2.5.7).

\subsection{Waktu Penerapan \textit{Post-Training Quantization} dan \textit{Quantization-Aware Training}}

Kuantisasi pascapelatihan (\textit{post-training quantization}, PTQ) menentukan skala dan \textit{zero point} dari model yang sudah dilatih, sehingga tidak memerlukan pelatihan ulang dan dapat diterapkan pada model hasil \textit{fine-tuning} \parencite{gholami2021survey}. Sebaliknya, \textit{quantization-aware training} (QAT) memasukkan operasi kuantisasi ke dalam proses pelatihan agar model beradaptasi dengan presisi rendah, dengan konsekuensi biaya komputasi karena pelatihan ulang \parencite{gholami2021survey}. Perbedaan keduanya dirangkum pada Tabel~\ref{tab:qat-ptq}.

\begin{table}[htbp]
  \centering
  \small
  \renewcommand{\arraystretch}{1.2}
  \caption{Perbandingan \textit{Quantization-Aware Training} dan \textit{Post-Training Quantization} \parencite{gholami2021survey, abdurrahman2023implementing}}
  \label{tab:qat-ptq}
  \begin{tabularx}{\textwidth}{@{}>{\raggedright\arraybackslash}p{3.2cm}>{\raggedright\arraybackslash}X>{\raggedright\arraybackslash}X@{}}
    \hline
    \textbf{Aspek} & \textbf{QAT} & \textbf{PTQ} \\
    \hline
    Waktu penerapan & Selama proses pelatihan & Setelah model selesai dilatih \\
    Pelatihan ulang & Diperlukan & Tidak diperlukan \\
    Biaya komputasi & Tinggi karena pelatihan ulang & Rendah, hanya penentuan parameter kuantisasi \\
    F1 rata-rata IndoBERT 4-bit & 74,94 \parencite{abdurrahman2023implementing} & 72,78 \parencite{abdurrahman2023implementing} \\
    Contoh penelitian & Q8BERT \parencite{zafrir2019q8bert} dan QAT pada IndoBERT \parencite{abdurrahman2023implementing} & PTQ pada IndoBERT \parencite{abdurrahman2023implementing} dan Mokey \parencite{zadeh2022mokey} \\
    \hline
  \end{tabularx}
\end{table}

\subsection{Granularitas \textit{Per-Tensor} dan \textit{Per-Channel}}

Granularitas menentukan cakupan bobot yang memakai parameter kuantisasi bersama \parencite{gholami2021survey}. Pada \textit{per-tensor} (\textit{coarse-grained}), satu skala dipakai untuk seluruh \textit{tensor}, sehingga $\alpha$ pada Persamaan~\ref{eq:skala} adalah nilai absolut terbesar seluruh \textit{tensor} $W$:
\begin{equation}
  s = \frac{\max(|W|)}{q_{\max}}
  \label{eq:skala-coarse}
\end{equation}
Pada \textit{per-channel} (\textit{fine-grained}), satu skala ditentukan untuk setiap baris keluaran $j$ dari \textit{tensor} bobot \parencite{gholami2021survey}:
\begin{equation}
  s_j = \frac{\max(|W_j|)}{q_{\max}}
  \label{eq:skala-fine}
\end{equation}
Kedua granularitas memakai rumus kuantisasi yang sama dan hanya berbeda pada cakupan nilai maksimum untuk menghitung skala. Sumbu $j$ adalah sumbu keluaran, yaitu baris pada matriks bobot lapisan linear, sedangkan pada \textit{embedding} satu baris adalah satu \textit{token}.

Contoh berikut memakai matriks bobot $2 \times 4$ pada INT4 dengan baris A (0,80; $-0{,}30$; 0,10; 0,55) dan baris B (0,20; 0,09; $-0{,}15$; 0,04). Pada \textit{per-tensor}, nilai absolut terbesar seluruh matriks adalah 0,80, sehingga satu skala $s = 0{,}80/7 \approx 0{,}1143$ dipakai untuk kedelapan bobot. Pada \textit{per-channel}, baris A tetap memakai $s_A = 0{,}80/7 \approx 0{,}1143$, sedangkan baris B memakai $s_B = 0{,}20/7 \approx 0{,}0286$ karena nilai terbesarnya 0,20. Setiap bobot dibagi dengan skala barisnya, dibulatkan, lalu dikalikan kembali dengan skala yang sama, dan hasilnya ditunjukkan pada Tabel~\ref{tab:contoh-granularitas}. Pada baris A, kedua granularitas menghasilkan nilai yang sama karena skalanya sama. Pada baris B, rata-rata \textit{error} turun dari 0,0321 pada \textit{per-tensor} menjadi 0,0057 pada \textit{per-channel}, atau 5,6 kali lebih kecil, karena bobot kecil pada baris B tidak lagi dibulatkan dengan langkah yang ditentukan oleh nilai 0,80 pada baris A.

\begin{table}[htbp]
  \centering
  \scriptsize
  \setlength{\tabcolsep}{3pt}
  \renewcommand{\arraystretch}{1.15}
  \caption{Contoh penghitungan \textit{per-tensor} dan \textit{per-channel} pada matriks bobot $2 \times 4$}
  \label{tab:contoh-granularitas}
  \begin{tabularx}{\textwidth}{@{}c>{\centering\arraybackslash}p{0.9cm}*{6}{>{\centering\arraybackslash}X}@{}}
    \hline
    \textbf{Baris} & \textbf{Bobot} & \textbf{\textit{Per-tensor} bilangan bulat} & \textbf{\textit{Per-tensor} dibaca kembali} & \textbf{\textit{Per-tensor} \textit{error}} & \textbf{\textit{Per-channel} bilangan bulat} & \textbf{\textit{Per-channel} dibaca kembali} & \textbf{\textit{Per-channel} \textit{error}} \\
    \hline
    A & 0,80 & 7 & 0,8000 & 0,0000 & 7 & 0,8000 & 0,0000 \\
    A & $-0{,}30$ & $-3$ & $-0{,}3429$ & 0,0429 & $-3$ & $-0{,}3429$ & 0,0429 \\
    A & 0,10 & 1 & 0,1143 & 0,0143 & 1 & 0,1143 & 0,0143 \\
    A & 0,55 & 5 & 0,5714 & 0,0214 & 5 & 0,5714 & 0,0214 \\
    B & 0,20 & 2 & 0,2286 & 0,0286 & 7 & 0,2000 & 0,0000 \\
    B & 0,09 & 1 & 0,1143 & 0,0243 & 3 & 0,0857 & 0,0043 \\
    B & $-0{,}15$ & $-1$ & $-0{,}1143$ & 0,0357 & $-5$ & $-0{,}1429$ & 0,0071 \\
    B & 0,04 & 0 & 0,0000 & 0,0400 & 1 & 0,0286 & 0,0114 \\
    \hline
  \end{tabularx}
\end{table}

Cakupan yang lebih kecil membuat nilai ekstrem pada satu bagian bobot tidak memengaruhi skala bagian lain, tetapi menambah jumlah skala yang harus disimpan \parencite{gholami2021survey}. Faktor skala ikut tersimpan di samping bobot terkuantisasi, sehingga ukuran berkas aktual lebih besar daripada ukuran teoretis. Tabel~\ref{tab:granularitas} menempatkan \textit{per-tensor} dan \textit{per-channel} di antara pilihan granularitas yang tersedia.

\begin{table}[htbp]
  \centering
  \small
  \renewcommand{\arraystretch}{1.2}
  \caption{Perbandingan tingkat granularitas kuantisasi \parencite{gholami2021survey}}
  \label{tab:granularitas}
  \begin{tabularx}{\textwidth}{@{}>{\raggedright\arraybackslash}p{3.4cm}>{\raggedright\arraybackslash}X>{\raggedright\arraybackslash}X>{\raggedright\arraybackslash}X@{}}
    \hline
    \textbf{Aspek} & \textbf{\textit{Per-tensor}} & \textbf{\textit{Per-channel}} & \textbf{\textit{Per-group}} \\
    \hline
    Cakupan satu skala & Satu \textit{tensor} & Satu baris keluaran & Satu kelompok elemen \\
    Jumlah skala pada model penelitian & 75 & 133.458 & Tidak diuji \\
    Variasi antarbaris & Tidak ditangani & Ditangani & Ditangani \\
    Variasi dalam satu baris & Tidak ditangani & Tidak ditangani & Ditangani \\
    \hline
  \end{tabularx}
\end{table}

\subsection{Penyimpanan Bobot INT4 dengan Teknik \textit{Bit-Packing}}

Memori umumnya dialamatkan per \textit{byte} (8 bit), sedangkan nilai INT4 hanya membutuhkan 4 bit. Tanpa penanganan khusus, satu nilai INT4 menempati satu \textit{byte} sehingga setengah kapasitas tidak terpakai. \textit{Bit-packing} memadatkan dua nilai INT4 ke dalam satu \textit{byte} melalui operasi \textit{bitwise} berupa penggeseran dan penggabungan bit, sehingga ukuran berkas bergantung pada cara penyimpanan dan bukan hanya pada jumlah bit per nilai.

Tanpa \textit{bit-packing}, ukuran berkas INT4 dua kali ukuran teoretisnya, yaitu 50\% dari ukuran FP16 dan bukan 25\%. Parameter yang tidak dikuantisasi dan faktor skala menambah kurang dari 0,5\% dari ukuran FP16. Pada kuantisasi tersimulasi, bobot dikembalikan ke presisi asal melalui \textit{dequantization} sebelum komputasi sehingga penghematan hanya terjadi pada penyimpanan, bukan pada waktu komputasi \parencite{gholami2021survey}. Karena itu, latensi tidak dapat diklaim membaik hanya karena bobot disimpan pada presisi lebih rendah (Subbab 2.6.2).

\section{Metrik Evaluasi dan Pengujian Signifikansi}

\subsection{Metrik \textit{Exact Match} dan F1}

\textit{Exact Match} (EM) dan F1 adalah metrik standar evaluasi SQuAD \parencite{rajpurkar2016squad}. Sebelum dibandingkan, jawaban prediksi dan jawaban acuan dinormalisasi dengan mengubah huruf menjadi huruf kecil, membuang tanda baca dan kata sandang bahasa Inggris (\textit{a}, \textit{an}, \textit{the}), serta menyeragamkan spasi \parencite{rajpurkar2016squad}. Pada SQuAD-id, pembuangan kata sandang bahasa Inggris hanya berdampak pada \textit{token} bahasa Inggris yang tidak diterjemahkan (seperti nama atau judul) dan berlaku setara pada kedua jawaban. Sebuah sampel bernilai 1 apabila jawaban prediksi cocok dengan salah satu jawaban acuan setelah normalisasi dan 0 apabila tidak, dan EM adalah proporsi sampel yang bernilai 1:
\begin{equation}
  \mathrm{EM} = \frac{\text{jumlah sampel yang cocok}}{N}
  \label{eq:em}
\end{equation}
\noindent dengan $N$ jumlah sampel.

F1 mengukur \textit{token overlap} antara jawaban prediksi dan jawaban acuan. \textit{Precision} ($P$) adalah proporsi \textit{token} prediksi yang juga muncul pada jawaban acuan, dan \textit{recall} ($R$) adalah proporsi \textit{token} jawaban acuan yang muncul pada prediksi \parencite{rajpurkar2016squad}. F1 dihitung dari keduanya:
\begin{equation}
  F1 = \frac{2 \cdot P \cdot R}{P + R}
  \label{eq:f1}
\end{equation}
Pada setiap sampel dipakai jawaban acuan yang menghasilkan F1 tertinggi, lalu F1 akhir adalah rata-rata seluruh sampel.

\subsection{Metrik Ukuran Model dan Latensi Inferensi}

Latensi inferensi adalah waktu yang diperlukan model untuk menghasilkan prediksi bagi satu \textit{batch} masukan \parencite{gholami2021survey}. Latensi bergantung pada perangkat, implementasi \textit{kernel}, dan cara penyimpanan bobot, sehingga dilaporkan sebagai metrik deskriptif, tidak diuji signifikansinya, dan tidak dipakai sebagai dasar klaim percepatan. Prosedur pengukurannya dijelaskan pada Subbab 3.2.7.

Efisiensi penyimpanan diukur melalui ukuran model dalam megabita (MB, 1 MB = $10^{6}$ byte), yaitu ukuran berkas setelah status model disimpan. Ukuran berkas aktual berbeda dari ukuran teoretis pada Subbab 2.4 karena faktor skala disimpan di samping bobot terkuantisasi, \textit{bias}, normalisasi lapisan, dan \textit{answer head} tidak dikuantisasi, dan \textit{bit-packing} tidak diterapkan (Subbab 2.5.7). Karena itu, ukuran teoretis dan ukuran berkas dilaporkan sebagai dua nilai yang berbeda.

\subsection{Uji Signifikansi dengan McNemar dan \textit{Paired Bootstrap}}

Selisih EM atau F1 antarkonfigurasi belum tentu mencerminkan perbedaan kinerja yang sesungguhnya, karena nilai tersebut dihitung pada sampel uji yang terbatas \parencite{dror2018hitchhiker}. Uji signifikansi diperlukan agar selisih yang teramati dapat dibedakan dari variasi acak pada sampel uji \parencite{dror2018hitchhiker}. Karena kedua model dievaluasi pada sampel uji yang sama, digunakan pengujian berpasangan.

Untuk EM yang bernilai biner pada setiap sampel, selisih dua model diuji dengan uji McNemar \parencite{dror2018hitchhiker, mcnemar1947}. Sampel dikelompokkan menurut kebenaran prediksi kedua model, dengan $b$ jumlah sampel yang benar pada model pertama tetapi salah pada model kedua dan $c$ sebaliknya. Statistik ujinya dengan koreksi kontinuitas \parencite{edwards1948continuity} adalah:
\begin{equation}
  \chi^2 = \frac{\left(|b - c| - 1\right)^2}{b + c}
  \label{eq:mcnemar}
\end{equation}
Apabila $b + c$ kurang dari 25, digunakan uji binomial eksak \parencite{fagerland2013mcnemar}.

Untuk F1 yang bernilai kontinu, \textit{paired bootstrap} mengambil sampel ulang dengan pengembalian dari himpunan sampel uji sebanyak $B$ kali, lalu menghitung selisih rata-rata skor kedua model pada setiap sampel ulang \parencite{bergkirkpatrick2012, efron1993}:
\begin{equation}
  \Delta^{*} = \bar{f}^{*}_{A} - \bar{f}^{*}_{B}
  \label{eq:bootstrap}
\end{equation}
\noindent dengan $\bar{f}^{*}_{A}$ dan $\bar{f}^{*}_{B}$ rata-rata skor model A dan model B pada satu sampel ulang. Penelitian ini memakai $B = 10.000$, sehingga diperoleh 10.000 nilai $\Delta^{*}$. Selisih dianggap signifikan apabila interval persentil 95\% dari $\Delta^{*}$ tidak memuat nol dengan taraf signifikansi $\alpha = 0{,}05$. Nilai $p$ dua sisi diperoleh dari dua kali proporsi sampel ulang yang selisihnya berada di sisi nol yang lebih kecil \parencite{efron1993}.

Perbandingan beberapa konfigurasi pada \textit{dataset} yang sama menambah jumlah uji berpasangan dan peluang kesalahan tipe I, sehingga nilai $p$ dikoreksi dengan metode Holm \parencite{holm1979}. Koreksi diterapkan terpisah untuk EM dan F1. $m$ nilai $p$ diurutkan dari terkecil, dan hipotesis ke-$k$ ditolak apabila $p_{(k)} \le \alpha/(m-k+1)$ serta seluruh hipotesis sebelumnya telah ditolak. Ukuran model bersifat deterministik sehingga tidak diuji, dan latensi diperlakukan sebagai metrik deskriptif.

\section{Penelitian Terdahulu dan Celah Penelitian}

\subsection{Ringkasan Penelitian Terdahulu}

Penelitian terdahulu pada Tabel~\ref{tab:penelitian-terdahulu} dipilih dengan empat syarat. Pertama, penelitian menerapkan kuantisasi tanpa pelatihan ulang atau pada model \textit{encoder} \textit{Transformer}. Kedua, penelitian dievaluasi pada GLUE atau SQuAD. Ketiga, penelitian terbit dalam lima tahun terakhir (2021--2026). Keempat, penelitian telah disitasi minimal lima kali menurut OpenAlex per 21 Oktober 2026. Urutannya dari yang paling relevan, yaitu PTQ \textit{weight-only} 4-bit pada BERT, hingga studi yang hanya berbagi masalah \textit{outlier} atau pengelompokan skala. SQuAD (termasuk SQuAD v1.1 dan v2.0) dipakai pada empat dari enam studi di tabel \parencite{zadeh2022mokey, bai2022towards, guo2023olive, wei2022outlier}. \textcite{gholami2021survey} adalah survei yang menjadi dasar teori Subbab 2.5 dan tidak dimasukkan ke tabel karena tidak melaporkan eksperimen.

\textcite{zadeh2022mokey} mengubah bobot BERT-Large menjadi indeks 4-bit tanpa \textit{fine-tuning} dan melaporkan F1 SQuAD v1.1 sebesar 93,17 dibanding 93,15 pada presisi penuh. \textcite{bai2022towards} meminimalkan \textit{error} rekonstruksi per modul, sedangkan \textcite{wei2022outlier} dan \textcite{guo2023olive} menangani \textit{outlier} aktivasi pada presisi 6-bit dan 4-bit. \textcite{bondarenko2021understanding} menganalisis kegagalan kuantisasi naif pada \textit{Transformer} dan memakai pengelompokan skala \textit{per-embedding-group}, sedangkan \textcite{kim2021ibert} membangun inferensi bilangan bulat penuh dengan \textit{speedup} 2,4 sampai 4,0 kali pada GPU T4.

Keenam penelitian tersebut dievaluasi pada bahasa Inggris. Penelitian pada bahasa Indonesia yang relevan terbit pada 2023 sampai 2026 dan belum mencapai lima sitasi pada tanggal tersebut, sehingga tidak dimasukkan ke tabel tetapi dibahas di sini. \textcite{abdurrahman2023implementing} menerapkan PTQ dan QAT pada IndoBERT dengan bobot 16, 8, 6, dan 4 bit untuk tujuh tugas IndoNLU, dan melaporkan F1 rata-rata PTQ 4-bit sebesar 72,78 dibanding 77,35 pada FP32, tanpa menyebut granularitas. Ramaputra dan Joddy (\citeyear{ramaputra2026evaluating}) menerapkan PTQ \textit{weight-only} dengan AWQ dan GPTQ pada model Qwen2.5 dan Gemma-3 untuk SQuAD-id, dengan EM Qwen2.5-14B 4-bit turun dari 48,61 menjadi 44,73 dan 43,80, sedangkan Qwen2.5-7B 4-bit berubah dari 42,05 menjadi 42,15 dan 42,34. \textcite{suhendra2026matrik} membandingkan model BERT berbahasa Indonesia pada SQuAD-id tanpa kuantisasi dan memakai \textit{bootstrap}, dengan F1 63,59 untuk \texttt{indobert-\allowbreak base-\allowbreak p1} dan 71,58 untuk \texttt{indolem/\allowbreak indobert-\allowbreak base-\allowbreak uncased}. \textcite{sanjaya2026quantization} membandingkan PTQ dan QAT pada IndoBERT untuk analisis sentimen tiga kelas dengan presisi FP16, INT8, dan INT4, dan melaporkan akurasi yang berubah tidak lebih dari 0,4 poin persentase.

{\setlength{\tabcolsep}{3pt}
\captionsetup[longtable]{font=normalsize}
\begingroup\scriptsize
\renewcommand{\arraystretch}{1.15}
\begin{longtable}{@{}c>{\raggedright\arraybackslash}p{1.55cm}>{\raggedright\arraybackslash}p{1.95cm}>{\raggedright\arraybackslash}p{2.45cm}>{\raggedright\arraybackslash}p{1.95cm}>{\raggedright\arraybackslash}p{1.25cm}>{\raggedright\arraybackslash}p{2.55cm}@{}}
  \caption{Penelitian terdahulu}
  \label{tab:penelitian-terdahulu} \\
  \hline
  \textbf{No} & \textbf{Referensi} & \textbf{Tujuan atau masalah} & \textbf{Model dan \textit{dataset}} & \textbf{Teknik} & \textbf{Metrik} & \textbf{Hasil utama} \\
  \hline
  \endfirsthead
  \hline
  \multicolumn{7}{c}{\tablename~\thetable{} (lanjutan)} \\
  \hline
  \textbf{No} & \textbf{Referensi} & \textbf{Tujuan atau masalah} & \textbf{Model dan \textit{dataset}} & \textbf{Teknik} & \textbf{Metrik} & \textbf{Hasil utama} \\
  \hline
  \endhead
  \hline
  \endlastfoot
  1 & Zadeh et al. (2022) & Metode \textit{fixed-point} biasanya menuntut pelatihan ulang, sehingga perlu cara agar model \textit{floating point} dapat dijalankan tanpa \textit{fine-tuning} & BERT-Base/Large, RoBERTa-Large, DeBERTa-XL; MNLI, STS-B, SQuAD v1.1 & PTQ \textit{weight-only} 4-bit: bobot menjadi indeks 4-bit ke \textit{dictionary} \textit{centroid} 16-bit & F1, akurasi & BERT-Large SQuAD: F1 93,15 (FP) $\rightarrow$ 93,17 (\textit{weight-only}), 92,22 bila aktivasi ikut dikuantisasi \\
  2 & Bai et al. (2022) & PTQ naif menurunkan akurasi model bahasa pralatih, sehingga diperlukan optimasi per modul & Model bahasa pralatih (keluarga BERT); GLUE dan SQuAD & PTQ per modul (MREM): \textit{error} rekonstruksi diminimalkan pada tiap modul & EM, F1, metrik GLUE & Hasil mendekati QAT pada GLUE dan SQuAD dengan biaya pelatihan lebih rendah \\
  3 & Guo et al. (2023) & \textit{Outlier} pada aktivasi membuat kuantisasi 4-bit bobot dan aktivasi gagal & BERT-Base/Large, BART-Base; GLUE (8 tugas), SQuAD v1.1 dan v2.0 & PTQ 4-bit bobot dan aktivasi (W4A4) dengan \textit{outlier-victim pair} & Akurasi, F1 & Penurunan akurasi kurang dari 1\% pada 4-bit, mengungguli PTQ 6-bit \textit{outlier suppression} \\
  4 & Wei et al. (2022) & \textit{Outlier} pada aktivasi dan \textit{embedding} menghambat kuantisasi presisi rendah & BERT, RoBERTa, BART; GLUE, SQuAD v1.1 dan v2.0, peringkasan & PTQ 6-bit: \textit{gamma migration} dan \textit{token-wise clipping} & EM, F1, metrik GLUE & PTQ BERT 6-bit mencapai level presisi penuh. Pada SQuAD v1.1 unggul 4,73\% (BERT) dan 15,55\% (RoBERTa) dibanding metode PTQ 6-bit terbaik sebelumnya \\
  5 & Bondarenko et al. (2021) & Kuantisasi naif gagal pada \textit{Transformer} & BERT-Base; GLUE & PTQ dan QAT dengan \textit{per-embedding-group quantization} & Metrik GLUE & GLUE 83,06 (FP32), 71,03 (W8A8), 70,25 (W32A8), 83,23 (W8A32) \\
  6 & Kim et al. (2021) & Operasi non-linear seperti GELU, Softmax, dan LayerNorm menghalangi kuantisasi \textit{integer-only} & BERT, RoBERTa; GLUE & Kuantisasi \textit{integer-only} INT8 (W8A8) dengan aproksimasi GELU, Softmax, dan LayerNorm & Metrik GLUE & Akurasi setara atau sedikit lebih tinggi dari FP32, dengan \textit{speedup} 2,4--4,0$\times$ pada GPU T4 \\
\end{longtable}
\endgroup
}

\noindent Keterangan: MREM (\textit{module-wise reconstruction error minimization}), GLUE (\textit{General Language Understanding Evaluation}), MNLI (\textit{Multi-Genre Natural Language Inference}), STS-B (\textit{Semantic Textual Similarity Benchmark}). Angka antarbaris tidak dapat dibandingkan langsung karena metrik dilaporkan pada \textit{dataset} dan tugas yang berbeda. Tabel ini menunjukkan selisih kinerja terhadap presisi penuh pada masing-masing penelitian, bukan peringkat metode.

\subsection{Celah Penelitian dan Kerangka Berpikir}

Tabel~\ref{tab:penelitian-terdahulu} dan penelitian berbahasa Indonesia pada Subbab 2.7.1 menunjukkan tiga celah. Pertama, kuantisasi bobot 4-bit pada model \textit{encoder} dilaporkan pada BERT berbahasa Inggris \parencite{zadeh2022mokey}. Pada IndoBERT, kuantisasi 4-bit baru dilaporkan sebagai F1 rata-rata tujuh tugas IndoNLU, termasuk FacQA yang berupa \textit{span extraction}, tanpa menyebut granularitas \parencite{abdurrahman2023implementing}. Pada SQuAD-id, kuantisasi baru dilaporkan untuk model \textit{decoder-only} \parencite{ramaputra2026evaluating}, dan studi pada model \textit{encoder} tidak menerapkan kuantisasi \parencite{suhendra2026matrik}. Kedua, pengelompokan skala pada \textit{encoder} dikaji pada BERT berbahasa Inggris dengan tugas GLUE \parencite{bondarenko2021understanding}. Granularitas relevan pada INT4 karena skema simetris hanya menyediakan 15 tingkat nilai ($-7$ sampai 7) dan satu nilai ekstrem dapat memperlebar langkah kuantisasi seluruh bobot dalam satu \textit{tensor} (Persamaan~\ref{eq:skala-coarse}), sedangkan \textit{per-channel} membatasi pengaruh itu pada satu baris (Persamaan~\ref{eq:skala-fine}). Ketiga, studi yang membandingkan model pada SQuAD-id memakai \textit{bootstrap} \parencite{suhendra2026matrik}, tetapi tidak untuk menguji selisih model terkuantisasi terhadap \textit{baseline} dengan koreksi perbandingan berganda. Dalam penelusuran penulis, ketiga hal tersebut belum ditemukan.

Penelitian ini menjawab celah tersebut dengan mengevaluasi \textit{weight-only} PTQ INT4 pada IndoBERT untuk SQuAD-id. Dua granularitas, \textit{per-tensor} dan \textit{per-channel}, dibandingkan dengan model FP16 menggunakan EM, F1, ukuran model, dan latensi, dan selisih EM dan F1 diuji signifikansinya. Alur penalaran ini dirangkum pada Gambar~\ref{fig:kerangka-berpikir}.

\begin{figure}[htbp]
  \centering
  \includegraphics[width=\textwidth,height=0.34\textheight,keepaspectratio]{"resources/Kerangka berpikir penelitian.png"}
  \caption{Kerangka berpikir penelitian}
  \label{fig:kerangka-berpikir}
\end{figure}

    %--------------------------------------------------------------------%
%
% Bab III — Metode Penelitian
%
%--------------------------------------------------------------------%

\chapter{Metode Penelitian}

\section{Jenis dan Alur Penelitian}

\subsection{Jenis Penelitian Eksperimental Kuantitatif}

Penelitian ini merupakan penelitian eksperimental kuantitatif. Perlakuan yang diberikan adalah kuantisasi \textit{weight-only} INT4 pada model yang telah di-\textit{fine-tune}, dengan dua granularitas sebagai variasi perlakuan, yaitu \textit{per-tensor} dan \textit{per-channel}. Kondisi pengujian dikendalikan agar identik, yaitu data uji, model dasar, lingkungan komputasi, dan prosedur evaluasi yang sama untuk seluruh konfigurasi.

Pendekatan eksperimental dipilih karena tiga alasan. Pertama, rumusan masalah menuntut pengukuran akibat suatu perlakuan, bukan pengamatan keadaan yang sudah ada. Kedua, perlakuan dapat diberikan secara terkendali karena kuantisasi hanya mengubah presisi bobot, sedangkan arsitektur model tidak berubah. Ketiga, hasil perlakuan dapat diukur secara kuantitatif melalui metrik baku tugas \textit{extractive question answering} (Subbab 2.6).

Penelitian ini tidak termasuk penelitian pengembangan perangkat lunak. Keluaran penelitian berupa bukti empiris mengenai pengaruh kuantisasi INT4 dan granularitasnya, bukan sistem atau aplikasi baru. Kode program kuantisasi disertakan sebagai lampiran agar hasil dapat direproduksi, dan kedudukannya adalah alat penelitian.

\subsection{Alur Penelitian}

Penelitian dimulai dengan studi literatur untuk menyusun rumusan masalah pada Bab I dan celah penelitian pada Subbab 2.7.2, lalu menetapkan tujuan dan batasan penelitian. Eksperimen dirancang dengan menetapkan data, model, konfigurasi kuantisasi, metrik, dan uji statistik. Data disiapkan dan model di-\textit{fine-tune} pada FP32 untuk tiga nilai \textit{seed}.

Model hasil \textit{fine-tuning} kemudian dijadikan \textit{baseline} dan model terkuantisasi, lalu ukurannya diukur. Ketiga konfigurasi diuji pada data uji, dianalisis secara deskriptif, dan diuji signifikansinya. Kesimpulan ditarik dan laporan disusun. Alur ini ditunjukkan pada Gambar~\ref{fig:alur-penelitian}.

\begin{figure}[htbp]
  \centering
  \includegraphics[width=0.90\textwidth,height=0.40\textheight,keepaspectratio]{"resources/Alur penelitian.png"}
  \caption{Alur penelitian}
  \label{fig:alur-penelitian}
\end{figure}

\section{Desain Eksperimen}

\subsection{Gambaran Umum Desain Eksperimen}

Desain eksperimen bertumpu pada satu model dasar yang sama untuk seluruh konfigurasi. Model IndoBERT-base fase 2 di-\textit{fine-tune} pada data latih dalam FP32, lalu bobot hasil \textit{fine-tuning} disalin tiga kali. Salinan pertama dikonversi ke FP16 sebagai \textit{baseline}. Salinan kedua dikuantisasi INT4 dengan granularitas \textit{per-tensor}, dan salinan ketiga dikuantisasi INT4 dengan granularitas \textit{per-channel}.

Pemisahan dilakukan pada kuantisasi, bukan pada pelatihan. Dengan cara ini ketiga konfigurasi berbagi bobot awal yang identik, sehingga perbedaan hasil hanya berasal dari presisi dan granularitas bobot. Ketiga konfigurasi dievaluasi pada data uji yang sama dengan prosedur yang sama, dan alurnya ditunjukkan pada Gambar~\ref{fig:desain-eksperimen}.

\begin{figure}[htbp]
  \centering
  \includegraphics[width=0.90\textwidth,height=0.40\textheight,keepaspectratio]{"resources/desain eksperimen.png"}
  \caption{Desain eksperimen untuk tiga konfigurasi}
  \label{fig:desain-eksperimen}
\end{figure}

\subsection{Lingkungan Eksperimen}

Seluruh eksperimen dijalankan pada lingkungan komputasi yang sama, mulai dari penyiapan data, \textit{fine-tuning}, kuantisasi, hingga pengukuran. Lingkungan tersebut dirangkum pada Tabel~\ref{tab:lingkungan}. Spesifikasi dicantumkan karena pengukuran latensi dan ukuran berkas bergantung pada perangkat keras, sistem operasi, dan versi pustaka.

\begin{table}[!ht]
  \centering
  \small
  \renewcommand{\arraystretch}{1.2}
  \caption{Lingkungan eksperimen penelitian}
  \label{tab:lingkungan}
  \begin{tabularx}{\textwidth}{@{}>{\raggedright\arraybackslash}p{4.2cm}X@{}}
    \hline
    \textbf{Komponen} & \textbf{Spesifikasi} \\
    \hline
    Platform & Kaggle Notebooks \\
    GPU & NVIDIA Tesla T4, 16 GB VRAM, satu GPU dipakai \\
    CPU & 4 vCPU \\
    Memori & 29 GB RAM \\
    Penyimpanan & 20 GB, ditambah \textit{dataset} yang diunggah permanen \\
    Batas eksekusi & 9 jam per sesi, kuota 30 jam GPU per minggu \\
    Bahasa pemrograman & Python \\
    Pustaka utama & PyTorch, \textit{transformers}, \textit{datasets}, NumPy \\
    Versi pustaka & Mengikuti \textit{image} Kaggle saat eksekusi dan dicatat pada lampiran \\
    \hline
  \end{tabularx}
\end{table}

Dua keputusan lingkungan perlu dijelaskan. Pertama, hanya satu GPU yang dipakai walaupun platform menyediakan dua, karena model 123,9 juta parameter yang dilatih pada 64.054 baris dapat ditampung pada satu GPU 16 GB sehingga pelatihan terdistribusi tidak diperlukan. Kedua, seluruh \textit{run} dijalankan pada jenis GPU yang sama, karena latensi yang diukur pada perangkat berbeda tidak dapat dibandingkan walaupun model dan data uji identik.

Kaggle dipilih karena menyediakan GPU T4 dengan kuota 30 jam per minggu dan batas sesi 9 jam. \textit{Dataset} diunggah sebagai \textit{dataset} permanen agar tidak diunduh ulang pada setiap sesi. Lokasi penelitian adalah lingkungan daring Kaggle Notebooks, sehingga penelitian tidak terikat pada satu ruang fisik.

\subsection{Data Penelitian dan Analisis Data Eksploratif}

Penelitian ini menggunakan \textit{dataset} SQuAD-id, yaitu terjemahan mesin SQuAD v2.0 ke bahasa Indonesia \parencite{muis2020sequence}. Karakteristik berkas mentah dirangkum pada Tabel~\ref{tab:karakteristik-dataset} di Subbab 2.3. Data bersifat sekunder dan tersedia untuk umum. Lisensi \textit{dataset} belum ditetapkan pada kartu \textit{dataset}, sehingga ketentuan penggunaan diperiksa sebelum data didistribusikan ulang.

Berkas data diunduh dari sumber publik yang tercantum pada kartu \textit{dataset}, lalu diunggah ke Kaggle sebagai \textit{dataset} permanen (Subbab 3.2.2). Instrumen pengumpulan data adalah \textit{dataset} itu sendiri, tanpa kuesioner atau wawancara. Setiap baris berisi satu konteks, satu pertanyaan, teks jawaban beserta posisi awalnya, dan penanda apakah pertanyaan tersebut tidak dapat dijawab.

Analisis data eksploratif dilakukan pada berkas mentah sebelum penyaringan. Pemeriksaan mencakup lima hal, yaitu keutuhan berkas dan jumlah baris tiap pembagian, jumlah artikel dan konteks unik, distribusi penanda pertanyaan tanpa jawaban, kesesuaian posisi jawaban dengan teksnya, serta panjang konteks dan pertanyaan dalam kata dan \textit{token}. Pemeriksaan posisi jawaban paling menentukan karena berkas mentah memuat 25.462 dari 86.821 baris terjawab pada data latih (29,3\%) yang posisi jawabannya tidak menunjuk pada teks jawaban.

Hasil pemeriksaan menunjukkan dua hal yang memengaruhi rancangan eksperimen. Pertama, satu konteks dipakai rata-rata 6,8 pertanyaan (130.319 baris pada 19.029 konteks unik), sehingga pemisahan data secara acak per baris akan menempatkan konteks yang sama pada data latih dan data validasi. Kedua, panjang konteks memiliki median 98 kata dan persentil ke-95 sebesar 192 kata, sehingga batas panjang masukan dapat ditetapkan dari distribusi tersebut (Subbab 3.2.4).

Pembagian data dilakukan pada tingkat artikel, bukan pada tingkat baris. Data uji memakai seluruh pembagian \textit{dev} asli yang memuat 35 artikel, dan tidak ada artikel yang muncul pada data latih sekaligus data uji. Data latih mentah memuat 441 artikel, dan setelah penyaringan dua artikel tidak menyisakan baris yang dapat dipakai sehingga tersisa 439 artikel.

Data validasi dibentuk dengan memisahkan 10\% dari 439 artikel tersebut, yaitu 44 artikel, menggunakan \textit{seed} tetap, dan 395 artikel sisanya menjadi data latih. Dari 70.051 baris hasil penyaringan, 5.997 baris masuk ke data validasi sehingga tersisa 64.054 baris untuk pelatihan. Pemisahan pada tingkat artikel membuat data validasi meniru kondisi data uji sekaligus mencegah kebocoran konteks, dan hasil pembagian akhir ditunjukkan pada Tabel~\ref{tab:pembagian-data}.

\begin{table}[!ht]
  \centering
  \small
  \renewcommand{\arraystretch}{1.2}
  \caption{Pembagian data setelah praproses}
  \label{tab:pembagian-data}
  \begin{tabularx}{\textwidth}{@{}Xrrr@{}}
    \hline
    \textbf{Ukuran} & \textbf{Latih} & \textbf{Validasi} & \textbf{Uji} \\
    \hline
    Jumlah baris & 64.054 & 5.997 & 5.244 \\
    Artikel & 395 & 44 & 35 \\
    Konteks unik & 17.136 & 1.574 & 1.200 \\
    Jumlah jawaban acuan & 64.054 & 5.997 & 16.428 \\
    Median panjang konteks (\textit{token}) & 135 & 136 & 137 \\
    \hline
  \end{tabularx}
\end{table}

\subsection{Praproses Data}

Praproses membuang baris pertanyaan tanpa jawaban, memperbaiki posisi jawaban, membuang duplikat, dan menyiapkan fitur masukan. Baris bertanda pertanyaan tanpa jawaban dibuang, yaitu 43.498 baris data latih dan 5.945 baris data uji, karena penelitian ini tidak menguji kemampuan model memutuskan bahwa suatu pertanyaan tidak dapat dijawab (Subbab 1.5).

Posisi jawaban diperbaiki pada baris yang ditandai tidak valid. Untuk setiap jawaban, konteks dibersihkan dari karakter tak terlihat dan spasi ganda, lalu pencocokan dilakukan bertingkat, yaitu \textit{exact matching}, \textit{exact matching} setelah normalisasi kapitalisasi dan diakritik, dan \textit{fuzzy matching} berbasis Levenshtein dengan ambang kemiripan 85. Hasil pencocokan dirapatkan ke batas kata, dan teks jawaban disimpan sebagai potongan konteks agar posisi dan isinya konsisten.

Perbaikan berhasil pada 8.757 baris data latih, terdiri atas 3.168 baris yang hanya memerlukan normalisasi dan 5.589 baris yang memerlukan \textit{fuzzy matching}. Sebanyak 16.705 baris data latih dan 676 baris data uji dibuang karena jawabannya tidak ditemukan pada konteks. Penulis memeriksa secara manual 30 jawaban hasil \textit{fuzzy matching} dan menemukan 29 jawaban benar (96,7\%), dengan interval kepercayaan Wilson 95\% sebesar 83\% sampai 99\%. Sampel ini kecil, sehingga angka tersebut dilaporkan sebagai gambaran kualitas perbaikan dan bukan sebagai jaminan.

Duplikat berdasarkan pasangan konteks dan pertanyaan dibuang, yaitu 65 baris data latih dan 8 baris data uji. Setiap contoh kemudian dibentuk dari pasangan pertanyaan dan konteks dengan panjang maksimum 384 \textit{token} dan \textit{stride} 128 \textit{token} antarjendela, sehingga konteks yang lebih panjang dipecah menjadi lebih dari satu jendela yang saling tumpang tindih. Panjang jawaban maksimum ditetapkan 30 \textit{token}.

Batas 384 \textit{token} ditetapkan dari distribusi panjang konteks hasil tokenisasi, yaitu median 135 \textit{token}, persentil ke-95 sebesar 261 \textit{token}, dan persentil ke-99 sebesar 349 \textit{token}. Hanya 0,70\% konteks yang melebihi 384 \textit{token} dan dipecah menjadi lebih dari satu jendela, dan seluruh jawaban berada di dalam setidaknya satu jendela sehingga tidak ada label yang hilang. Setiap contoh menyimpan tiga masukan, yaitu \textit{input\_ids}, \textit{attention\_mask}, dan \textit{token\_type\_ids}, serta dua label berupa posisi awal dan posisi akhir jawaban.

\subsection{Model IndoBERT-base dan Pengaturan \textit{Fine-Tuning}}

Model yang dipakai adalah \texttt{indobert-base-p2} dengan 12 lapisan \textit{encoder}, ukuran \textit{hidden} 768, 12 \textit{head} perhatian, ukuran lapisan \textit{feed-forward} 3.072, dan kosakata 50.000 \textit{token}. Model untuk tugas \textit{question answering} tidak memiliki lapisan \textit{pooler}, sehingga jumlah parameternya 123.852.290. Model ini dipralatih pada korpus Indo4B \parencite{wilie2020indonlu}. Fase 2 dipilih karena melanjutkan pralatih pada panjang urutan maksimum 512 \textit{token}, sedangkan fase 1 hanya 128 \textit{token}, dan penelitian ini memakai jendela hingga 384 \textit{token} (Subbab 2.2.2).

\textit{Task head} ditambahkan di atas keluaran \textit{encoder} untuk memprediksi posisi awal dan posisi akhir jawaban. \textit{Head} ini terdiri atas satu lapisan linear berukuran 768 ke 2, yaitu 1.538 parameter (0,001\% dari seluruh parameter). Keluaran model berupa skor untuk setiap posisi \textit{token}. \textit{Span} dengan jumlah skor awal dan akhir tertinggi, dengan posisi awal tidak melewati posisi akhir dan panjang paling banyak 30 \textit{token}, dipilih sebagai jawaban, dan untuk pertanyaan yang konteksnya terbagi menjadi lebih dari satu jendela dipakai \textit{span} dengan skor tertinggi di antara seluruh jendela.

\textit{Fine-tuning} dilakukan dengan pengaturan pada Tabel~\ref{tab:hyperparameter}. \textit{Learning rate} dan ukuran \textit{batch} berada pada rentang yang direkomendasikan \textcite{devlin2019bert} untuk \textit{fine-tuning} BERT, sedangkan jumlah \textit{epoch} maksimum 6 dibatasi oleh \textit{early stopping}. Pada setiap \textit{epoch}, model dievaluasi pada data validasi, dan bobot dengan F1 validasi tertinggi disimpan sebagai \textit{checkpoint}. Fungsi kerugian adalah \textit{cross-entropy} pada posisi awal dan posisi akhir jawaban, dengan nilai kerugian keduanya dirata-ratakan. Seluruh pengaturan ditetapkan sebelum pelatihan dan tidak disetel menggunakan data uji.

\begin{table}[!ht]
  \centering
  \small
  \renewcommand{\arraystretch}{1.2}
  \caption{Pengaturan \textit{fine-tuning}}
  \label{tab:hyperparameter}
  \begin{tabularx}{\textwidth}{@{}>{\raggedright\arraybackslash}p{6.2cm}X@{}}
    \hline
    \textbf{Parameter} & \textbf{Nilai} \\
    \hline
    Model dasar & \texttt{indobert-base-p2} \\
    Panjang urutan maksimum & 384 \textit{token} \\
    \textit{Stride} antarjendela & 128 \textit{token} \\
    Panjang jawaban maksimum & 30 \textit{token} \\
    Ukuran \textit{batch} & 16 \\
    Fungsi kerugian & \textit{Cross-entropy} posisi awal dan akhir, dirata-ratakan \\
    \textit{Optimizer} & AdamW \\
    \textit{Learning rate} & $2 \times 10^{-5}$ \\
    Jumlah \textit{epoch} maksimum & 6 \\
    \textit{Early stopping} & \textit{patience} 2 pada F1 validasi, ambang perbaikan 0,001 \\
    Pemilihan \textit{checkpoint} & F1 validasi tertinggi, bobot terbaik dipulihkan \\
    \textit{Seed} & 42, 1337, dan 2024 \\
    \hline
  \end{tabularx}
\end{table}

Tiga \textit{seed} dipakai untuk mengukur variasi hasil akibat proses pelatihan, mengikuti praktik pada penelitian kuantisasi model berbahasa Indonesia yang melaporkan tiga \textit{seed} \parencite{sanjaya2026quantization}. Setiap \textit{seed} menghasilkan satu \textit{checkpoint} yang menjadi dasar bagi ketiga konfigurasi.

\subsection{Implementasi Kuantisasi \textit{Weight-Only} INT4}

Kuantisasi diterapkan pada bobot model setelah \textit{fine-tuning} selesai, tanpa pelatihan ulang dan tanpa data kalibrasi. Skema yang dipakai adalah kuantisasi seragam simetris INT4 dengan bilangan bulat pada rentang $-7$ sampai 7 (Persamaan~\ref{eq:simetris} dan Persamaan~\ref{eq:skala} pada Subbab 2.5.3). Aktivasi tidak dikuantisasi, sehingga konfigurasi ini termasuk \textit{weight-only} PTQ (Subbab 2.5.4). Nilai $\alpha$ diambil dari bobot itu sendiri, sehingga tidak ada hiperparameter kuantisasi yang disetel menggunakan data uji.

Model hasil \textit{fine-tuning} FP32 disalin agar bobot asli tetap tersedia, lalu seluruh lapisan linear dan \textit{embedding} dikumpulkan dengan menelusuri modul secara rekursif. Untuk setiap \textit{tensor} bobot dihitung nilai absolut terbesar $\alpha$, yaitu satu $\alpha$ per \textit{tensor} pada konfigurasi \textit{per-tensor} (Persamaan~\ref{eq:skala-coarse}) dan satu $\alpha$ per baris pada konfigurasi \textit{per-channel} (Persamaan~\ref{eq:skala-fine}), lalu skala dihitung sebagai $s = \alpha / 7$. Setiap bobot dibagi dengan skalanya, dibulatkan ke bilangan bulat terdekat, dan dibatasi pada $-7$ sampai 7. Modul asli kemudian diganti dengan modul terkuantisasi yang menyimpan bilangan bulat dalam \textit{container} 8-bit beserta skala dalam FP32. Contoh penghitungan dengan angka tercantum pada Tabel~\ref{tab:contoh-format} dan Tabel~\ref{tab:contoh-granularitas}.

Skala dihitung dari bobot FP32 hasil \textit{fine-tuning}, bukan dari bobot FP16, agar pembulatan tidak terjadi dua kali. Pada inferensi, bilangan bulat dikalikan dengan skalanya melalui \textit{dequantization} sehingga kembali menjadi bobot FP16 sebelum operasi linear dijalankan, dan aktivasi tetap pada FP16. Karena itu, penghematan berada pada ukuran penyimpanan, bukan pada waktu komputasi. Bobot INT4 disimpan dalam \textit{container} 8-bit tanpa \textit{bit-packing}, sehingga ukuran bobot terkuantisasi dua kali ukuran teoretisnya (Subbab 2.5.7).

Prosedur yang sama diterapkan pada \textit{checkpoint} ketiga \textit{seed}, karena prosedur ini hanya beroperasi pada \textit{tensor} bobot dan tidak bergantung pada data. Perbedaan antar\textit{seed} hanya berasal dari bobot hasil \textit{fine-tuning}.

Kuantisasi mencakup seluruh matriks bobot, termasuk \textit{embedding}, karena \textit{embedding} menyumbang 31,3\% dari seluruh parameter dan melewatkannya akan membuat 31,3\% parameter tetap berukuran FP16 sehingga penghematan ukuran yang menjadi tujuan penelitian berkurang. \textit{Bias}, normalisasi lapisan, dan \textit{answer head} tidak dikuantisasi dan disimpan dalam FP16, sama seperti pada \textit{baseline}. Sumbu \textit{per-channel} adalah baris \textit{tensor}, yaitu satu keluaran pada lapisan linear dan satu \textit{token} pada \textit{embedding} kata, posisi, dan tipe \textit{token}. Modul yang dikuantisasi, jumlah parameter, dan jumlah skalanya ditunjukkan pada Tabel~\ref{tab:modul-kuantisasi}.

\begin{table}[!ht]
  \centering
  \small
  \renewcommand{\arraystretch}{1.2}
  \caption{Modul IndoBERT-base untuk \textit{question answering} dan jumlah skala kuantisasinya}
  \label{tab:modul-kuantisasi}
  \begin{tabularx}{\textwidth}{@{}>{\raggedright\arraybackslash}Xrrrr@{}}
    \hline
    \textbf{Kelompok modul} & \textbf{Jumlah \textit{tensor}} & \textbf{Jumlah parameter} & \textbf{Dikuantisasi} & \textbf{Skala \textit{per-channel}} \\
    \hline
    \textit{Embedding} kata, posisi, dan tipe \textit{token} & 3 & 38.794.752 & ya & 50.514 \\
    Proyeksi atensi (\textit{query}, \textit{key}, \textit{value}, keluaran) & 48 & 28.311.552 & ya & 36.864 \\
    Lapisan \textit{feed-forward} (12 lapisan) & 24 & 56.623.104 & ya & 46.080 \\
    \textit{Bias} dan normalisasi lapisan & -- & 121.346 & tidak & -- \\
    \textit{Answer head} & -- & 1.536 & tidak & -- \\
    \hline
    Total & 75 & 123.852.290 & 123.729.408 & 133.458 \\
    \hline
  \end{tabularx}
\end{table}

Skala \textit{per-channel} menambah 0,2\% dari ukuran FP16, sedangkan skala \textit{per-tensor} hanya menambah 300 \textit{byte}. Ukuran model hasil hitungan ditunjukkan pada Tabel~\ref{tab:ukuran-model}, dengan 1 MB sama dengan $10^{6}$ \textit{byte}.

\begin{table}[!ht]
  \centering
  \small
  \renewcommand{\arraystretch}{1.2}
  \caption{Ukuran model hasil hitungan dengan parameter tak terkuantisasi FP16 dan skala FP32}
  \label{tab:ukuran-model}
  \begin{tabularx}{\textwidth}{@{}Xrr@{}}
    \hline
    \textbf{Konfigurasi} & \textbf{Ukuran (MB)} & \textbf{Rasio terhadap FP16} \\
    \hline
    FP16 & 247,7 & 100\% \\
    INT4 \textit{per-tensor} dalam \textit{container} 8-bit & 124,0 & 50,1\% \\
    INT4 \textit{per-channel} dalam \textit{container} 8-bit & 124,5 & 50,3\% \\
    \hline
  \end{tabularx}
\end{table}

Ukuran berkas aktual diukur setelah model disimpan dan dibandingkan dengan hitungan pada Tabel~\ref{tab:ukuran-model}.

\subsection{Skenario dan Pengukuran Pengujian}

Pengujian dilakukan terhadap tiga konfigurasi pada data uji yang sama, dengan tiga \textit{seed} sebagai ulangan. Setiap konfigurasi diukur dengan empat metrik, yaitu EM, F1, ukuran model, dan latensi inferensi. Definisi keempat metrik dibahas pada Subbab 2.6, sedangkan subbab ini menetapkan skenario dan prosedur pengukurannya.

EM dan F1 diukur pada setiap \textit{seed}, sehingga terdapat sembilan nilai untuk masing-masing metrik dari tiga konfigurasi dan tiga \textit{seed}. Ukuran model diukur sekali untuk setiap konfigurasi karena nilainya tidak bergantung pada \textit{seed}. Latensi diukur pada setiap konfigurasi menggunakan model \textit{seed} 42 dengan satu \textit{batch} berisi 16 contoh berpanjang 384 \textit{token} yang diambil dari data uji. Pengukuran diawali 10 iterasi \textit{warm-up} dan dilanjutkan 100 pengulangan, dengan sinkronisasi CUDA sebelum dan sesudah setiap pengulangan, lalu dilaporkan median dan persentil ke-95 waktu per \textit{batch}. Rangkuman skenario ditunjukkan pada Tabel~\ref{tab:skenario}.

\begin{table}[!ht]
  \centering
  \small
  \renewcommand{\arraystretch}{1.2}
  \caption{Skenario pengujian}
  \label{tab:skenario}
  \begin{tabularx}{\textwidth}{@{}c>{\raggedright\arraybackslash}p{2.4cm}>{\raggedright\arraybackslash}p{2.4cm}>{\raggedright\arraybackslash}p{2.8cm}>{\raggedright\arraybackslash}p{1.6cm}>{\raggedright\arraybackslash}X@{}}
    \hline
    \textbf{No} & \textbf{Konfigurasi} & \textbf{Presisi bobot} & \textbf{Granularitas} & \textbf{Ulangan} & \textbf{Metrik} \\
    \hline
    1 & \textit{Baseline} & FP16 & -- & 3 \textit{seed} & EM, F1, ukuran, latensi \\
    2 & Kuantisasi & INT4 & \textit{per-tensor} & 3 \textit{seed} & EM, F1, ukuran, latensi \\
    3 & Kuantisasi & INT4 & \textit{per-channel} & 3 \textit{seed} & EM, F1, ukuran, latensi \\
    \hline
  \end{tabularx}
\end{table}

Data uji berisi 5.244 pertanyaan dengan 16.428 jawaban acuan. Evaluasi EM dan F1 memakai jawaban acuan dengan skor tertinggi untuk setiap pertanyaan, sesuai prosedur standar SQuAD \parencite{rajpurkar2016squad}, dan normalisasi teks mengikuti Subbab 2.6.1.

\section{Teknik Analisis Hasil Penelitian}

\subsection{Analisis Deskriptif}

Analisis deskriptif menggambarkan selisih kinerja dan efisiensi antara konfigurasi kuantisasi dan \textit{baseline}. Untuk EM dan F1 dilaporkan rata-rata dan simpangan baku sampel (pembagi $n-1$) dari tiga \textit{seed}, serta selisih terhadap \textit{baseline} pada \textit{seed} yang sama. Untuk ukuran model dilaporkan ukuran berkas aktual dalam megabita dan rasionya terhadap FP16. Untuk latensi dilaporkan median dan persentil ke-95 tanpa uji statistik.

Pertukaran antara kinerja dan efisiensi dianalisis dengan menyandingkan penurunan F1 (dalam poin) dan pengurangan ukuran model (dalam persen) pada setiap granularitas. Dari penyandingan ini terlihat granularitas mana yang memberi penghematan ukuran terbesar untuk penurunan F1 terkecil.

Selisih EM dan F1 dihitung berpasangan per \textit{seed} agar variasi akibat pelatihan tidak tercampur dengan pengaruh kuantisasi. Selisih antargranularitas kemudian dibandingkan dengan simpangan baku antar\textit{seed}. Bila selisih antargranularitas lebih kecil daripada simpangan baku tersebut, kedua granularitas dilaporkan tidak dapat dibedakan secara deskriptif. Karena hanya ada tiga \textit{seed}, simpangan baku dipakai sebagai gambaran distribusi dan bukan sebagai dasar inferensi.

\subsection{Uji Signifikansi Statistik}

Analisis deskriptif menunjukkan besarnya selisih, tetapi tidak menunjukkan apakah selisih tersebut berasal dari perlakuan atau dari variasi acak pada sampel uji. Karena itu dilakukan uji signifikansi berpasangan, dengan pemilihan uji sesuai sifat metrik pada Tabel~\ref{tab:uji-metrik}.

\begin{table}[!ht]
  \centering
  \small
  \renewcommand{\arraystretch}{1.2}
  \caption{Pemilihan uji signifikansi per metrik}
  \label{tab:uji-metrik}
  \begin{tabularx}{\textwidth}{@{}l>{\raggedright\arraybackslash}p{5.6cm}X@{}}
    \hline
    \textbf{Metrik} & \textbf{Sifat nilai} & \textbf{Uji} \\
    \hline
    EM & biner per sampel, berpasangan & uji McNemar dengan koreksi kontinuitas \\
    F1 & kontinu per sampel, berpasangan & \textit{paired bootstrap} 10.000 sampel ulang \\
    Ukuran model & deterministik & tidak diuji \\
    Latensi & bergantung perangkat dan implementasi & tidak diuji, dilaporkan deskriptif \\
    \hline
  \end{tabularx}
\end{table}

Untuk EM, hasil kedua konfigurasi disusun menjadi tabel kontingensi, dengan $b$ jumlah pertanyaan yang benar pada konfigurasi pertama tetapi salah pada konfigurasi kedua dan $c$ sebaliknya. Statistik uji dihitung dengan Persamaan~\ref{eq:mcnemar}. Hipotesis nol menyatakan bahwa kedua konfigurasi memiliki proporsi jawaban benar yang sama ($b = c$ dalam harapan), sedangkan hipotesis alternatif menyatakan proporsi itu berbeda. Bila jumlah sampel yang berbeda ($b + c$) kurang dari 25, dipakai uji binomial eksak \parencite{fagerland2013mcnemar}. Untuk F1, selisih rata-rata skor dihitung pada setiap sampel ulang dengan Persamaan~\ref{eq:bootstrap}. Sampel ulang diambil pada tingkat pertanyaan (5.244 pertanyaan), dengan indeks yang sama untuk kedua konfigurasi. Hipotesis nol menyatakan bahwa rata-rata selisih F1 kedua konfigurasi sama dengan nol, sedangkan hipotesis alternatif menyatakan selisih itu tidak nol. Nilai $p$ dua sisi dan interval persentil 95\% dari $\Delta^{*}$ dihitung dari 10.000 sampel ulang tersebut \parencite{bergkirkpatrick2012, efron1993}, dan selisih dilaporkan bersama intervalnya.

Ukuran model tidak diuji karena nilainya dihitung secara deterministik dari jumlah parameter dan tipe datanya, tanpa proses pengukuran yang mengandung variasi acak. Latensi tidak diuji karena nilainya bergantung pada perangkat dan implementasi, dan hanya diukur pada satu model dengan pengulangan pada perangkat yang sama.

Terdapat tiga perbandingan berpasangan, yaitu \textit{baseline} terhadap \textit{per-tensor}, \textit{baseline} terhadap \textit{per-channel}, dan \textit{per-tensor} terhadap \textit{per-channel}. Ketiga perbandingan ini meningkatkan peluang kesalahan tipe I, sehingga nilai $p$ dikoreksi dengan metode Holm \parencite{holm1979}. Koreksi diterapkan terpisah untuk setiap metrik dengan $m = 3$, dan dijalankan di dalam masing-masing \textit{seed}.

Suatu perbedaan dinyatakan signifikan bila signifikan setelah koreksi Holm pada taraf $\alpha = 0{,}05$ di ketiga \textit{seed} dengan arah yang sama. Syarat ketiga \textit{seed} ini lebih ketat daripada satu uji tunggal, sehingga koreksi tambahan antar\textit{seed} tidak diperlukan. Uji berpasangan pada setiap \textit{seed} hanya menangkap variasi sampel uji untuk model yang telah dilatih, sedangkan variasi akibat \textit{seed} pelatihan dilaporkan secara deskriptif. Karena itu, kesimpulan penelitian ini tidak mencakup seluruh kemungkinan hasil pelatihan.

\section{Jadwal Penelitian}

Penelitian direncanakan berlangsung selama 6 bulan, yaitu Juli sampai Desember 2026, dimulai dari studi literatur dan penyusunan proposal, dilanjutkan penyiapan data, pelatihan, kuantisasi, pengujian, dan penyusunan laporan. Rincian kegiatan per bulan ditunjukkan pada Tabel~\ref{tab:jadwal}.

\begin{table}[!ht]
  \centering
  \small
  \renewcommand{\arraystretch}{1.2}
  \caption{Jadwal penelitian}
  \label{tab:jadwal}
  \begin{tabularx}{\textwidth}{@{}c>{\raggedright\arraybackslash}X*{6}{c}@{}}
    \hline
    \textbf{No} & \textbf{Kegiatan} & \textbf{Jul} & \textbf{Agu} & \textbf{Sep} & \textbf{Okt} & \textbf{Nov} & \textbf{Des} \\
    \hline
    1 & Studi literatur dan penyusunan proposal & \checkmark & \checkmark & & & & \\
    2 & Penyusunan kajian pustaka dan kerangka berpikir & & \checkmark & \checkmark & & & \\
    3 & Analisis data eksploratif dan praproses data & & & \checkmark & \checkmark & & \\
    4 & Implementasi kode program pelatihan dan kuantisasi & & & & \checkmark & & \\
    5 & \textit{Fine-tuning} tiga \textit{seed} dan pemilihan \textit{checkpoint} & & & & \checkmark & \checkmark & \\
    6 & Penerapan kuantisasi dan pengukuran ukuran model & & & & & \checkmark & \\
    7 & Pengujian, analisis deskriptif, dan uji signifikansi & & & & & \checkmark & \\
    8 & Penyusunan Bab IV dan Bab V & & & & & \checkmark & \checkmark \\
    9 & Pemeriksaan akhir dan penyusunan laporan & & & & & & \checkmark \\
    \hline
  \end{tabularx}
\end{table}

    %--------------------------------------------------------------------%
%
% Bab IV — Evaluasi dan Pembahasan
%
% @author Petra Novandi
%
%--------------------------------------------------------------------%
%
% Berisi tujuan pengujian, skenario, hasil pengujian, dan pembahasan.
% Ganti teks placeholder dengan isi sebenarnya.
%
%--------------------------------------------------------------------%

\chapter{Evaluasi dan Pembahasan}

Bab ini mempresentasikan pengujian solusi yang dirancang pada Bab
III dan membahas apakah tujuan serta rumusan masalah telah
terpenuhi. Hasil yang ditampilkan harus dapat dipertanggungjawabkan
secara objektif.

\section{Tujuan Pengujian}

Tuliskan tujuan pengujian yang spesifik dan terukur. Setiap tujuan
sebaiknya berkaitan langsung dengan tujuan tugas akhir dan
rumusan masalah di Bab I.

\blindtext%

\section{Skenario Pengujian}

Jelaskan metode, lingkungan, data uji, dan langkah-langkah
pengujian yang dilakukan. Skenario harus cukup jelas agar pengujian
dapat direplikasi oleh pembaca.

\blindtext%

\section{Hasil Pengujian}

Sajikan hasil pengujian secara faktual menggunakan tabel, grafik,
atau angka. Hindari interpretasi panjang di subbab ini. Fokuskan
pada data dan temuan yang diperoleh.

    \subsection{Subbab}

    \blindtext%

    \subsubsection{Subsubbab}

    \blindtext%

\blindtext%

\section{Pembahasan}

Interpretasikan hasil pengujian dan kaitkan dengan rumusan masalah
serta tujuan tugas akhir. Jelaskan apakah solusi yang diusulkan
berhasil, keterbatasannya, dan perbandingannya dengan studi terkait
jika relevan.

\blindtext%

    %--------------------------------------------------------------------%
%
% Bab V — Penutup
%
% @author Petra Novandi
%
%--------------------------------------------------------------------%
%
% Berisi kesimpulan dan saran. Ganti teks placeholder dengan isi
% sebenarnya.
%
%--------------------------------------------------------------------%

\chapter{Penutup}

Bab ini merangkum hasil tugas akhir dan memberikan saran untuk
pengembangan lebih lanjut. Kesimpulan harus menjawab rumusan masalah
dan tujuan yang telah dirumuskan di awal laporan.

\section{Kesimpulan}

Tuliskan kesimpulan utama berdasarkan hasil analisis, perancangan,
dan evaluasi. Setiap poin kesimpulan sebaiknya merujuk pada tujuan
atau rumusan masalah yang telah diselesaikan.

\blindtext%

\section{Saran}

Berikan saran yang konstruktif untuk pengembangan solusi, perbaikan
metodologi, atau penelitian lanjutan. Saran dapat mencakup aspek
teknis, evaluasi, maupun implementasi di lingkungan nyata.

\blindtext%


    % Daftar Rujukan (gaya APA). heading=bibintoc membuat judul sekaligus
    % entri Daftar Isi pada halaman pertama daftar rujukan.
    {\setlength{\emergencystretch}{2em}\printbibliography[heading=bibintoc, title={DAFTAR RUJUKAN}]}

    %----------------------------------------------------------------%
    % Lampiran — label "Lampiran <angka>".
    %----------------------------------------------------------------%
    \appendix
    \renewcommand{\thechapter}{\arabic{chapter}}
    % Entri Daftar Isi untuk lampiran: "Lampiran 1", bukan "BAB I".
    \renewcommand{\tocchapapp}{Lampiran}
    \renewcommand{\tocchapnum}{\arabic{chapter}}
    \titleformat{\chapter}[display]
      {\normalsize\bfseries\centering}
      {Lampiran\ \thechapter}{0pt}
      {\MakeUppercase}

    % Entri "LAMPIRAN" dihapus: sebelumnya menunjuk halaman terakhir Daftar
    % Rujukan, sedangkan Lampiran 1 dan 2 sudah punya entri sendiri.

    %--------------------------------------------------------------------%
%
% Lampiran A — Instrumen Pengujian
%
% @author Petra Novandi
%
%--------------------------------------------------------------------%
%
% Berisi instrumen atau alat bantu pengujian. Ganti teks placeholder
% dengan isi sebenarnya.
%
%--------------------------------------------------------------------%

\chapter{Instrumen Pengujian}

\blindtext%

    %--------------------------------------------------------------------%
%
% Lampiran B — Rincian Kasus Uji
%
% @author Petra Novandi
%
%--------------------------------------------------------------------%
%
% Berisi rincian kasus uji. Ganti teks placeholder dengan isi
% sebenarnya.
%
%--------------------------------------------------------------------%

\chapter{Rincian Kasus Uji}

\blindtext%


    %--------------------------------------------------------------------%
%
% Riwayat Hidup (Pedoman UM Lampiran 13) — naratif, spasi tunggal,
% tanpa foto. Batasi pendidikan sampai jenjang di atas SMA.
%
%--------------------------------------------------------------------%

\clearpage
\addcontentsline{toc}{chapter}{RIWAYAT HIDUP}

\begin{center}
    \bfseries RIWAYAT HIDUP
\end{center}

\begin{singlespace}
\blindtext%
\end{singlespace}
     % Riwayat Hidup

\end{document}
